Textos de Emma Goldman

Vítimas da moralidade (1913)

Os aspectos sociais do controle de natalidade (1916)

Novamente sobre o movimento de controle de natalidade (1916)

O camaleônico sufrágio feminino (1917)

Louise Michel (1923)

Mulheres heroicas da revolução russa (1925)

As visões de Emma sobre o amor (1926)

A luta do feminismo não foi em vão (1926)

O elemento sexual da vida (s/ data)




\chapter[Vítimas da moralidade (1913)]{Vítimas da moralidade\footnote[*]{Texto originalmente publicado na revista
  anarquista ``\emph{Mother Earth}'' (1913).}}

Não faz muito tempo que participei de um encontro organizado por Anthony
Comstock, que há quarenta anos tem sido o guardião da moral americana.
Nunca ouvi antes de qualquer palanque, uma divagação mais incoerente e
ignorante.

A questão que se apresentou para mim, ao ouvir a conversa lugar-comum e
fanática desse homem, foi: ``Como é possível que alguém tão limitado e
sem inteligência exerça o poder de censor e ditador sobre uma nação
supostamente democrática?'' É verdade que Comstock tem a lei para
respaldá-lo. Quarenta anos atrás, quando o puritanismo, era inclusive
mais desenfreado do que hoje, capaz de apagar completamente qualquer luz
de razão e progresso, Comstock conseguiu, através de maquinações
obscuras e conspiração política, fazer passar uma lei que lhe deu
controle absoluto do Departamento responsável pelas agências dos
correios {[}\emph{Post Office Departament}{]} -- um controle que se
mostrou desastroso à liberdade de imprensa, assim como ao direito de
privacidade do cidadão americano.\footnote{Mesmo antes da promulgação da
  lei de 1873 que ficou conhecida pelo seu nome, Anthony Comstock, já
  atuava na condição de agente da \emph{YMCA's Society for the
  Suppression of Vice} da cidade de Nova York -- cuja tarefa era
  censurar o material ``obsceno'' enviado pelos correios e,
  eventualmente, prender os seus remetentes e destinatários. Numa
  determinada ocasião, ainda no seu primeiro ano de trabalho em 1872,
  Comstock chegou a prender, de um lance, quarenta homens. Em realidade,
  a lei de 1873 é uma correção às ``brechas'' supostamente deixadas pela
  lei de 1872 (por sua vez uma revisão de uma lei federal de 1865) que
  justamente tinha por objetivo de interditar o comércio e postagem,
  mesmo se de caráter privado de material considerado ``obsceno''. Para
  Comstock, a lei de 1872 era insuficiente, porque não incluía sob a
  insígnia de ``obsceno'', propagandas de livros de conteúdo subversivo,
  além de não determinar de modo claro os meios através dos quais
  deveriam ser executadas as penalidades. Com a promulgação da lei,
  Comstock se tornou agente especial do \emph{Post Office Department},
  dotado de todo poder e instrumentos para determinar o que era ou não
  material obsceno e, por conseguinte, punir os envolvidos.}

Desde então, Comstock invadiu os aposentos privados das pessoas,
confiscou a sua correspondência pessoal, assim como trabalhos
artísticos, e estabeleceu um sistema de espionagem e suborno que
deixaria a Rússia envergonhada. Seja como for, as leis não dão conta de
explicar o poder de Anthony Comstock. Há algo além, ainda mais terrível
do que a lei. É a própria estreiteza do espírito puritano, tal como
representado nas mentes estéreis das \emph{Young-Men-and-Old-Maid's}
{[}Associações Cristãs de Moços e Solteironas{]}, \footnote{Aqui Goldman
  faz uma brincadeira maldosa, pois ao invés de se referir à
  \emph{Women's Christian Temperance Union} (que em tradução livre seria
  algo como União da Temperança Cristã das Mulheres) ela substitui a
  expressão ``Women'', pela expressão ``Old Maid'', isto é: solteirona.}
\emph{Christian Union}, \emph{Temperance Union}, \emph{Sabbath Union},
\emph{Purity League} etc. Um espírito que é absolutamente cego às mais
simples manifestações da vida; representante, portanto, de toda
estagnação e decadência. Tal como nos dias que precederam a guerra civil
americana {[}\emph{ante-bellum days}{]}, esses fósseis velhos lamentam a
terrível imoralidade do nosso tempo. Ciência, arte, literatura, drama
estão à mercê de uma censura intolerante munida de recursos legais, o
que tem como resultado que a América, apesar de todos as suas
declarações presunçosas em nome do progresso e da liberdade, ainda está
imersa no provincialismo mais denso.

As menores nações da Europa podem se gabar de possuir uma arte livre dos
grilhões da moralidade, uma arte que tem a coragem de retratar os
grandes problemas sociais de nosso tempo. Com a faca afiada da análise
crítica, corta cada úlcera social, cada erro, exigindo mudanças
fundamentais e a transvaloração dos valores aceitos.

Sátira, sagacidade, humor, assim como os modos de expressão mais sérios
e intensos, são empregados para pôr a nu as mentiras das nossas
convenções sociais e morais. Na América, procurar-se-ia em vão por tal
meio de expressão, mesmo porque qualquer tentativa nesse sentido está
interditada pelo rígido regime levado a cabo pelo ditador moral e sua
panelinha.

O que de mais próximo disso temos por aqui, são os nossos
\emph{muckrakers},\footnote{\emph{Muckraker} era a expressão utilizada
  para designar os membros de um grupo de jornalistas e escritores que
  se concentraram em elaborar matérias jornalísticas (e obras
  literárias) para denunciar a corrupção nas relações entre governo e
  grandes indústrias, além de outras instituições de prestígio. Boa
  parte das suas denúncias exerceu papel relevante (no sentido de
  pressionar) na elaboração e promulgação de projetos de lei, como, por
  exemplo, os que fortaleciam a proteção de trabalhadores e
  consumidores. Em geral, \emph{os muckrakers} são historicamente
  alocados na chamada Era Progressiva dos Estados Unidos que vai dos
  1890 até a entrada do país na Primeira Guerra Mundial, em 1917.} que,
sem sombra de dúvida, prestaram um ótimo serviço em termos econômicos e
sociais. Quer os \emph{muckrakers} tenham ou não ajudado a mudar
efetivamente as condições, ao menos eles arrancaram a máscara do rosto
mentiroso dessa nossa sociedade arrogante e convencida.

Infelizmente, a Mentira da Moralidade ainda persiste na sua perseguição,
e com toda a pompa, já que ninguém se atreve a mexer com o santo dos
santos. É certo, porém, que nenhuma outra superstição é tão prejudicial
ao crescimento e tão dotada do poder de debilitar e paralisar as mentes
e corações do povo, quanto a superstição da Moralidade.

O aspecto mais patético, e de certo modo desencorajador, dessa situação
toda é certo elemento presente nos liberais e mesmo nos radicais, enfim,
presente nos homens e mulheres aparentemente livres dos fantasmas
religiosos e sociais. Pois, ante o monstro da Moralidade, eles estão tão
prostrados quanto os mais devotos -- o que é uma prova adicional da
profundidade em que o verme da moralidade se introduziu no sistema de
suas vítimas e quão amplas e rigorosas devem ser as medidas para
expulsá-lo mais uma vez.

Não é preciso dizer que a sociedade é obcecada por mais de um tipo de
moral. É um fato que toda instituição da atualidade tem seu próprio
padrão moral. Seja como for, as instituições não poderiam se manter, não
fosse pela religião, que funciona como um escudo, e pela moralidade, que
funciona como uma máscara. Isso explica o interesse que os exploradores
ricos têm tanto na religião, quanto na moralidade. Os ricos pregam,
promovem e financiam ambas, como um investimento que gera excelentes
retornos. Por meio da religião, eles paralisaram o entendimento do povo,
assim como a moralidade escraviza o espírito. Em outras palavras, a
religião e a moralidade são um chicote muito mais eficiente, do que o
porrete e o revolver, para manter o povo submisso.

Para ilustrar: a Moralidade da Propriedade estabelece que essa
instituição é sagrada. Ai de quem se atrever a questionar a santidade da
propriedade ou pecar contra ela! Mesmo que todos saibam que a
Propriedade é um roubo; que representa o esforço acumulado de milhões de
pessoas que são desprovidas de propriedades. E, o que é mais terrível,
quanto mais assolada pela pobreza estiver a vítima da Moralidade da
Propriedade, maior respeito e temor ela terá por esse senhor. É, por
isso, que ouvimos de pessoas progressistas, até mesmo dos chamados
trabalhadores com consciência de classe, a condenação moral de métodos
como a sabotagem e a ação direta, justamente porque eles atingem a
Propriedade.

Na verdade, se as próprias vítimas estão tão cegas pela Moralidade da
Propriedade, o que se pode esperar dos seus senhores? Ao que parece, já
passou da hora de encarar o fato de que enquanto os trabalhadores não
tiverem perdido o respeito pelo instrumento que lhes escraviza
materialmente, a eles estão vedadas quaisquer esperanças.

\asterisc

O que, porém, me preocupa mais seriamente é o efeito da Moralidade sobre
as mulheres. Esse efeito é tão desastroso e tão paralisante que muitas
das mais progressistas dentre as minhas irmãs nunca o superaram
completamente.

É a Moralidade que condena a mulher à posição de celibatária, prostituta
ou ainda de chocadeira irresponsável e imprudente de um sem-número de
crianças infelizes.

Primeiramente, no que diz respeito ao celibato, ele torna o humano
semelhante a uma planta sedenta e murcha. Quando ainda uma flor jovem e
bela, ela se apaixona por um rapaz respeitável. Mas a Moralidade decreta
que, salvo o caso de que ele possa se casar com a garota, ela nunca deve
conhecer o arrebatamento do amor, o êxtase da paixão, que atinge o seu
cume na relação sexual. O jovem respeitável pode até estar disposto a se
casar, mas a Moralidade da Propriedade, da Família e todas as demais
Moralidades Sociais decretam que ele primeiro deve fazer fortuna,
economizar o suficiente para providenciar um lar e ser capaz de
sustentar a sua família. Os jovens são, assim, obrigados a esperar, não
raro, por muitos e longos cansativos anos.

Enquanto isso, o jovem respeitável, excitado com o relacionamento e
contato diário com sua namorada, procura uma saída para a natureza em
troca de dinheiro. Em noventa e nove por cento dos casos, ele será
infectado, de modo que quando estiver materialmente apto a se casar,
infectará a esposa e os possíveis filhos. E o que dizer da jovem flor,
com todas as suas fibras ardendo com o fogo da vida, com todo o seu ser
demandando amor e paixão? Ela não tem saída -- o que se manifesta em
dores de cabeça, insônia, histeria; torna-se amargurada, encrenqueira, e
logo se transforma num ser desbotado, murcho e sem alegria, um estorvo
para ela mesma e para todos os outros. Não é de admirar que Stirner
preferisse a \emph{grisette}\footnote{Expressão utilizada para designar
  jovens que complementavam a renda com a prostituição.} à donzela, cuja
virtude torna amarga e ressequida.

Não há nada mais patético, mais terrível, do que essa vítima embolorada
por uma Moralidade igualmente embolorada. Isso se aplica com força ainda
maior às jovens da classe média alta do que às do povo. Por conta da
necessidade econômica, as últimas são arremessadas na selva da vida
desde a mais tenra idade; elas crescem, lado a lado, com seus
companheiros homens seja na fábrica, na oficina, ou nas brincadeiras e
festas. O resultado é uma expressão mais normal dos seus instintos
físicos. Além disso, os rapazes e moças do povo não são moldados de modo
tão inflexível pelos fatores externos, e frequentemente se lançam ao
chamado do amor e da paixão, independentemente dos costumes e tradição.

Não obstante, a garota de classe média, nervosa e hipersexualizada,
protegida na estreiteza do confinamento das tradições familiares e
sociais, guardada por milhares de olhos, temerosa de sua própria sombra
-- tem de voltar o desejo do seu ser mais profundo, o desejo por um
homem ou por filhos, para gatos, cães, canários ou mesmo estudos
bíblicos. Tal é a cruel máxima da Moralidade, que diariamente exclui o
amor, a luz e a alegria da vida de inúmeras vítimas.

Agora, passemos à prostituta. Apesar das leis, decretos, perseguições e
prisões; apesar da segregação, dos registros, das cruzadas viciosas e
outros dispositivos semelhantes, a prostituta é o verdadeiro fantasma da
nossa época. Ela varre as planícies como se fosse um incêndio, atingindo
cada canto da vida, devastando, destruindo.

No final das contas, ela está retribuindo, em uma escala muito pequena,
a maldição e os horrores que a sociedade disseminou no seu caminho.
Exaurida pela perambulação ao longo eras, perseguida e obrigada a
peregrinar de um lado para o outro à mercê de todos, ela é, não
obstante, a Nêmesis dos tempos modernos, o anjo vingador a empunhar
impiedosamente a espada de fogo. Não tem ela o homem em seu poder? E,
através dele, o lar, a criança, a humanidade. É dessa maneira, portanto,
que ela mata, embora seja ela quem é assassinada do modo mais brutal. O
que foi que a criou? De onde ela vem? Da Moralidade, da moralidade que é
impiedosa nas suas atitudes para com as mulheres. Uma vez que a mulher
se atreva a ser ela mesma, a ser honesta para com a sua natureza, para
com a vida, não há retorno: ela é ejetada da jurisdição e proteção da
sociedade. A prostituta se torna vítima da Moralidade, assim como a
solteirona ressequida e murcha também é sua vítima. Mas a prostituta é
ainda vitimada por outras forças, principalmente pela Moralidade da
Propriedade, que obriga a mulher a se vender como mercadoria sexual por
um dólar a cada vez, quando fora do casamento, ou por quinze dólares por
semana, quando no solo sagrado do matrimônio. Essa última condição é sem
dúvida mais segura, mais respeitada, mais reconhecida, mas, entre as
duas formas de prostituição, a garota da rua é a menos hipócrita, a
menos degradada, já que o seu comércio não se vale da máscara piedosa da
hipocrisia; e, ainda assim, não obstante, é ela quem é perseguida,
extorquida, ultrajada e banida pelos mesmos poderes que a criaram: pelo
financista, pelo padre, pelo moralista, pelo juiz, pelo carcereiro e
pelo policial e, não podemos esquecer, por sua irmã bem salvaguardada,
respeitosamente virtuosa -- que é a mais implacável e brutal na sua
perseguição à prostituta.

A moralidade e sua vítima, a mãe -- que imagem terrível! Existe de fato
algo mais terrível, mais criminoso do que a glorificada e sagrada função
da maternidade? A mulher, física e mentalmente inapta para ser mãe, mas
condenada a procriar; a mulher, taxada economicamente em cada centelha
da sua energia, mas forçada a procriar; a mulher, amarrada a um homem
que detesta, cuja visão a enche de horror, mas moldada para procriar; a
mulher, cansada e exaurida pelo processo de gestação, mas coagida a
procriar, cada vez mais e mais. Que coisa hedionda, essa tão elogiada
maternidade! Não é de admirar que milhares de mulheres se arrisquem pela
via da mutilação, e prefiram até a morte a essa maldição cruelmente
imposta pelo fantasma da Moralidade. Cinco mil mulheres são anualmente
sacrificadas no altar desse monstro, que não tolera prevenção, mas que
supostamente seria uma cura para o aborto. Cinco mil guerreiras na
batalha pela sua liberdade física e espiritual, assim como o são muitos
milhares mais que preferiram se tornar aleijadas e mutiladas, ao invés
de dar à luz numa sociedade erguida sobre a decadência e destruição.

Será que é por não querer assumir responsabilidades, ou por não ter amor
aos seus filhos, que a mulher moderna é levada aos meios mais drásticos
e perigosos para evitar a gravidez? Apenas as pessoas dotadas de uma
mentalidade extremamente superficial e intolerante podem fazer tal
acusação. Não fosse isso, eles saberiam que a mulher moderna
simplesmente se tornou consciente da raça humana, sensível às
necessidades e direitos da criança, assim como à unidade da raça, e que,
portanto, a mulher moderna tem um senso de responsabilidade e
humanidade, que são completamente diferentes dos da sua avó.

Com a guerra econômica a destruir todo o seu entorno, com os conflitos,
a miséria, o crime, as doenças e a insanidade postas bem diante dos seus
olhos, com inúmeras crianças pequenas trituradas e destruídas, como
poderia uma mulher autoconsciente e consciente da própria humanidade vir
a se tornar mãe? A Moralidade não pode responder a essa pergunta. Só
pode doutrinar, coagir ou condenar -- e quantas mulheres são fortes o
suficiente para enfrentar essa condenação, para desafiar os dogmas
morais? Poucas, de fato. E, assim, terminam por preencher as fábricas,
os reformatórios, as casas de repouso, as prisões, os manicômios, ou
simplesmente morrem na tentativa de impedir uma gravidez indesejada. Oh,
a Maternidade, que crimes são cometidos em vosso nome! Que multidão é
colocada aos seus pés, Moralidade, destruidora da vida!

Felizmente, a Aurora está emergindo do caos e da escuridão. A mulher
está despertando, ela está se livrando do pesadelo da Moralidade; ela já
não ficará acorrentada. No seu amor pelo homem, ela não está mais
preocupada com o conteúdo de sua carteira, mas com a riqueza da sua
natureza, que por si só é a fonte da vida e da alegria. Tampouco ela
precisa da sanção do Estado. Seu amor é sanção suficiente para ela.
Assim, ela pode se entregar ao homem de sua escolha, como as flores
estão entregues ao orvalho e à luz, em liberdade, beleza e êxtase.

Através da sua consciência renascida como unidade, como personalidade,
como construtora da raça humana, ela se tornará mãe apenas se desejar o
filho e se puder dar à criança, mesmo antes de seu nascimento, tudo o
que sua natureza intelecto sejam capazes de alcançar: harmonia, saúde,
conforto, beleza e, acima de tudo, compreensão, reverência e amor, que é
o único solo verdadeiramente fértil para uma nova vida, para um novo
ser.

A Moralidade não pode aterrorizar quem se elevou além do bem e do mal. E
embora possa continuar devorando suas vítimas, essa mesma Moralidade é
completamente impotente diante do espírito moderno, que brilha em toda a
sua glória na testa do homem e da mulher libertos e sem medo.

\chapter[Os aspectos sociais do controle de natalidade (1916)]{Os aspectos sociais do controle de natalidade\footnote[*]{Texto
  originalmente publicado no jornal \emph{Mother Earth.} (1916)}}

Como já foi sugerido, para que um gênio seja criado, a natureza tem de
usar todos os seus recursos e demora uma centena de anos para levar a
cabo esta sua difícil tarefa. Se isto é verdade, a natureza leva ainda
mais tempo para criar uma grande ideia. Afinal, ao criar um gênio a
natureza se concentra em uma personalidade, enquanto uma ideia deve
eventualmente se tornar a herança de uma raça, devendo, por isso, ser
necessariamente mais difícil de moldar.

Faz apenas cento e cinquenta anos que um grande homem concebeu uma
grande ideia, Robert Thomas Malthus, o pai do controle de natalidade.
Que tenha demorado tanto tempo para a raça humana perceber a
grandiosidade desta ideia é apenas mais uma prova da morosidade da mente
humana. Não é possível entrar numa discussão detalhada acerca do mérito
da controvérsia colocada por Malthus, a saber, que a terra não é fértil
ou rica o suficiente para suprir as necessidades de uma raça
excessivamente abundante. Certamente, se olharmos por sobre as
trincheiras e os campos de batalha da Europa, descobriremos que, em
certa medida, a sua premissa estava correta. No entanto, tenho certeza
de que se Malthus vivesse hoje, ele concordaria com todos os estudiosos
da sociedade e revolucionários no sentido de que se as massas continuam
a ser pobres e os ricos tornam-se cada vez mais ricos, não é porque
esteja faltando fertilidade à terra e riqueza para suprir as
necessidades de uma raça abundante, mas porque a terra é monopolizada
nas mãos de poucos para a exclusão de muitos.

O capitalismo, que ainda estava de fraldas no tempo de Malthus, cresceu
desde então, tornando-se um enorme e insaciável monstro. Ele ruge
através do apito e da máquina -- ``Mandem suas crianças para mim, vou
torcer os seus ossos; vou extrair o seu sangue; vou roubá-las do seu
florescimento'' --, pois o capitalismo tem um apetite insaciável.

É através da sua maquinaria destrutiva que o militarismo e o capitalismo
proclamam: ``Mandem os seus filhos para mim, irei treiná-los e
discipliná-los até que toda a humanidade tenha sido programada neles;
até que eles se tornem autômatos prontos para atirar e matar sob o
comando dos seus mestres''. O capitalismo não pode se realizar sem o
militarismo e uma vez que as massas fornecem o material a ser destruído
nas trincheiras e campos de batalha, o capitalismo necessita de uma raça
numerosa.

Os chamados bons tempos são aqueles em que o capitalismo engole massas
de pessoas, para depois, novamente lançá-las no tempo da ``depressão
industrial''. Essa massa humana supérflua que incha as fileiras de
desempregados e que representa a maior ameaça nos tempos modernos, é
designada pelos nossos políticos e economistas burgueses de margem de
trabalho. Para eles, sob nenhuma circunstância deve a margem de trabalho
ser diminuída, caso contrário a sagrada instituição conhecida como
civilização capitalista será prejudicada. E assim, os economistas
políticos, em conjunto com todos financiadores do regime capitalista são
a favor de uma raça numerosa e excessiva e, portanto, opõem-se ao
controle de natalidade.

No entanto, a teoria de Malthus contém muito mais verdade do que ficção.
Dado o seu caráter ``moderno'', já não se baseia em especulação, mas em
outros fatores que estão relacionados e entrelaçados às profundas
mudanças sociais que estão acontecendo em todos os lugares.

Primeiro, há o aspecto científico, a contenda entre os mais eminentes
cientistas que nos dizem que um organismo sobrecarregado e subnutrido
não pode reproduzir prole saudável. Além dessa contenda dos cientistas,
estamos sendo confrontados com um fato terrível, agora inclusive
reconhecido pelas pessoas mais ignorantes, a saber: que uma procriação
indiscriminada e incessante da parte das massas sobrecarregadas e
subnutridas resultou no aumento de crianças deficientes, aleijadas e
infelizes. Este fato é tão alarmante que despertou os reformistas
sociais para a necessidade da criação de uma espécie de Câmara de
Compensação da mente que pudesse apurar as causas e efeitos do aumento
no número de crianças aleijadas, surdas, idiotas e cegas. Por sabermos
que os reformistas só aceitam a verdade quando ela se tornou óbvia até
para o indivíduo mais tapado da sociedade, já não há necessidade de
discutir os resultados de uma procriação indiscriminada.

Em segundo lugar, há o despertar intelectual da mulher que desempenha um
papel não pouco significativo em benefício do controle de natalidade. Ao
longo das eras, a mulher vem carregando o seu fardo. Cumpre o seu dever
mil vezes maior do que o de qualquer soldado em campo de batalha.
Afinal, o trabalho do soldado é tirar vidas. Para isso, ele é pago pelo
Estado, elogiado por charlatões políticos e apoiado pela histeria
pública. Já a função da mulher é dar à luz e, em relação a isso, nem os
políticos, nem a opinião pública demonstraram, em qualquer momento, a
mínima disposição de retribuir a vida que a mulher tem dado. Ao longo
das eras, ela se encontra de joelhos ante o altar do dever imposto por
Deus, pelo capitalismo, pelo Estado e pela moralidade. Na atualidade,
porém, ela está acordando desse seu sonho de longa data. Ela está se
libertando do pesadelo do passado; ela se virou em direção à luz que
anuncia, com voz clara, que ela não mais fará parte do crime de trazer
ao mundo crianças infelizes apenas para serem moídas e transformadas em
pó pelas rodas do capitalismo e para serem dilaceradas em pedaços nas
trincheiras e campos de batalha. E quem haverá de dizer a ela que não
seja assim? Afinal de contas, é a mulher quem arrisca a sua saúde e
sacrifica a sua juventude para a reprodução da raça. Certamente, ela
deveria estar em posição de decidir quantas crianças deve trazer ao
mundo, se elas devem ser trazidas ao mundo junto ao homem que ama e
porque ela deseja a criança, ou se devem nascer do ódio e da aversão.

Ademais, é consenso entre os médicos mais sérios que a reprodução
constante da parte das mulheres resulta no que em termos leigos é
chamado ``problemas femininos'': uma condição lucrativa para médicos
inescrupulosos. Mas por qual razão a mulher exaure o seu organismo numa
gravidez perpétua?

É precisamente por conta dessa razão que as mulheres devem ter o
conhecimento que as possibilite se recuperar por um período de três a
cinco anos entre cada gestação, o que por si só lhes daria bem-estar
físico e mental e a oportunidade de cuidar melhor das crianças já
existentes.

Mas não são apenas as mulheres que estão começando a perceber a
importância do controle de natalidade. Homens, especialmente da classe
trabalhadora, estão aprendendo a reconhecer que famílias numerosas são
como pedras amarradas em seus pescoços, deliberadamente impostas a eles
pelas forças reacionárias da sociedade; porque uma família numerosa
paralisa o cérebro e embota os músculos das massas de trabalhadores.
Nada amarra mais os trabalhadores à servidão do que uma ninhada de
crianças e é exatamente isso que os oponentes do controle de natalidade
querem.

Como os vencimentos de um homem com uma família numerosa são miseráveis,
ele não pode arriscar nem mesmo um pouco. Assim, permanece preso à
rotina, concedendo e se encolhendo diante de seu mestre, apenas para
ganhar sequer dá para alimentar as suas muitas boquinhas. Ele não se
atreve a participar de organizações revolucionárias; ele não se atreve a
fazer greve; ele não se atreve a expressar uma opinião. As massas de
trabalhadores, finalmente, acordaram para a necessidade do controle de
natalidade como meio de libertar a si mesmas do jugo terrível e ainda
mais como meio capaz de fazer algo por aqueles que já existem, ao evitar
que mais crianças venham ao mundo.

Por último, embora não menos importante, uma mudança na relação entre os
sexos, ainda que não abranja um número muito grande de pessoas, está se
fazendo sentir numa minoria considerável. Como no passado, para um
grande número de homens da atualidade, a mulher continua a ser mero
objeto, um meio para um fim; basicamente, um meio físico para um fim.
Mas há homens que querem mais do que isso de uma mulher; perceberam que
mesmo que todo homem se emancipe das superstições do passado, isso não
poderia mudar em nada a estrutura social caso a mulher não tome parte
com ele na grandiosa luta social. Lentamente, mas de modo firme, esses
homens têm aprendido que se a mulher desperdiça sua substância em
gestações eternas, resguardos e lavagens de fraldas, ela tem pouco tempo
para o que quer que seja. Muito menos, inclusive, para as questões que
absorvem e agitam o pai de seus filhos. Por exaustão física e estresse
nervoso, a mulher se torna um obstáculo no caminho do homem e
frequentemente seu inimigo mais amargo. É, portanto, para a sua própria
proteção e também em nome da companhia e amizade da mulher amada, que
muitos homens a querem livre da terrível imposição de reprodução
constante da vida, sendo, portanto, favoráveis ao controle de
natalidade.

Seja qual for o ângulo, a questão do controle de natalidade deve ser
considerada, é o assunto mais imperioso dos tempos modernos e, em sendo
assim, não pode ser censurado através de perseguições, prisões ou de uma
conspiração do silêncio. Aqueles que se opõem ao movimento do controle
de natalidade afirmam que o fazem em nome da maternidade. Todos os
charlatães políticos matraqueiam sobre essa tal maternidade maravilhosa,
que quando examinamos mais de perto, vemos que essa mesma maternidade
tem, ao longo dos séculos, dedicado os seus rebentos a Moloch. Além
disso, uma vez que as mães são obrigadas a trabalhar pesado, ao longo de
horas, para ajudar no sustento das criaturas que trouxeram ao mundo sem
assim o desejar, essa conversa sobre maternidade nada mais é do que
hipocrisia. Dez por cento das mulheres casadas da cidade de Nova York
têm de ajudar a ganhar o pão. A maioria delas ganha o salário bastante
profícuo de \$280 por ano. Quem ousa falar das belezas da maternidade
diante de tal crime?

Mas e se considerarmos mães com melhores salários, o que pensaremos
delas? Não faz muito tempo que nosso Conselho de Educação declarou que
não seria permitido a professoras que fossem mães continuar com o
magistério. Ainda que os antiquados cavaleiros em questão tenham sido
compelidos pela opinião pública a reconsiderar a sua decisão, é
absolutamente certo que quando uma professora se torna mãe, ano após
ano, ela perde a sua posição. Essa é a sina da mãe casada; e o que
acontece com a mãe solteira? Ou alguém tem dúvidas de que há milhares de
mães solteiras? Elas lotam as nossas lojas, fábricas e indústrias, estão
em todos os lugares, não por escolha, mas por necessidade econômica. Em
sua existência tediosa e monótona, a única alegria que resta é
provavelmente a da atração sexual, que sem os métodos de prevenção
invariavelmente conduz a abortos. Em segredo e com pressa, milhares de
mulheres são sacrificadas por causa de abortos realizados por médicos
charlatões e parteiras ignorantes. No entanto, os poetas e os políticos
louvam a maternidade. Não houve crime maior contra a mulher do que esse.

Nossos moralistas sabem disso, ainda que persistam na defesa da
reprodução indiscriminada de crianças. Eles nos dizem que limitar a
procriação é uma inclinação completamente moderna, porque a mulher
moderna está perdida moralmente e deseja fugir da responsabilidade. Em
resposta a isso, é necessário pontuar que a inclinação de controlar a
prole é tão velha quanto a raça humana. Temos como autoridade para essa
contenda o eminente médico alemão Dr. Theilhaber, que compilou dados
históricos para provar que esta inclinação era comum aos hebreus,
egípcios, persas e muitas tribos dos índios americanos. O medo da
criança era tão grande que muitas mulheres se valiam do mais hediondo
dos métodos para não trazer uma criança indesejada ao mundo. Dr.
Theilhaber enumera cinquenta e sete métodos.\footnote{Felix A.
  Theilhaber (1884 -1956) foi um médico e sexólogo judeu, militante da
  reforma sexual e da legalização do controle de natalidade e fundador
  de uma das primeiras clínicas de Berlim para o controle de natalidade.
  Nessa passagem, Goldman está se referindo à obra \emph{Das Sterile
  Berlin}, publicada em 1913.} Esses dados são de grande importância na
medida em que dissipam a superstição de que a mulher quer se tornar mãe
de uma família numerosa.

Não, não é que responsabilidade esteja faltando à mulher, mas é porque
ela tem muita responsabilidade que exige saber como prevenir a
concepção. Nunca na história do mundo, a mulher teve tanta consciência
da humanidade como ela tem hoje. Nunca antes ela esteve tão aberta a ver
na criança, não apenas na sua, mas em toda a criança, a unidade da
sociedade, o canal através do qual cada homem e cada mulher tem de
passar; o mais forte fator na construção de um novo mundo. É por esta
razão que o controle de natalidade se fundamenta em terreno sólido.

Foi-nos dito que se a lei, contida nos estatutos, tiver resultado de
discussões sobre como prevenir um determinado crime, os fatores de
prevenção já não precisam ser discutidos. Em resposta a isso, quero
dizer que não é o movimento do controle de natalidade que deve ser
descartado, mas a própria lei. Afinal, é para isso que as leis servem,
para serem feitas e desfeitas. Como eles ousam exigir que a vida se
submeta às leis? Devemos ficar atados às leis pelo resto das nossas
vidas, apenas porque algum fanático ignorante, na sua própria limitação
de mente e coração, foi bem-sucedido em aprovar uma lei na época em que
homens e mulheres eram prisioneiros de superstições morais e religiosas?
Eu compreendo perfeitamente por que juízes e carcereiros estão amarrados
às leis. Precisam dela para o seu sustento; em nome da sua função na
sociedade. Mas, às vezes, até juízes progridem. Chamo a atenção de
vocês, aqui, para a decisão a favor do controle de natalidade, do juiz
Gatens de Portland, Oregon. ``Parece-me que o problema com o nosso
pessoal atualmente é que há muito puritanismo. Ignorância e puritanismo
sempre foram a corda amarrada no pescoço do progresso. Todos nós sabemos
que muitas coisas estão erradas na nossa sociedade; que estamos sofrendo
de muitos males, mas nós não temos coragem para levantar e admitir, e
quando uma pessoa chama nossa atenção para algo que já sabemos, fingimos
decência e nos sentimos indignados.'' Este certamente é o problema com a
maioria dos nossos legisladores e com aqueles que se opõem ao controle
de natalidade.

Eu serei julgada nas \emph{Special Sessions} no dia 5 de
abril.\footnote{Como a legislatura e o processo penal dos Estados Unidos
  são específicos ao país, optou-se por manter aqui o nome no original.}
Não sei qual será o resultado, e não me importo.\footnote{Sob as leis de
  Comstock, fornecer informações sobre métodos contraceptivos era
  considerado crime, mesmo no caso de profissionais de saúde. A alegação
  era a de que esse tipo de informação promoveria a depravação sexual.
  Goldman foi presa em 11 de fevereiro daquele ano, poucos dias após
  falar para cerca de 3000 pessoas sobre métodos contraceptivos. No
  julgamento em questão, foi condenada e optou por passar quinze dias na
  prisão, ao invés de pagar a fiança de U\$100.} O medo de ir para
cadeia pelas próprias ideias é tão forte entre os intelectuais
americanos que é o que faz o movimento ser tão pálido e fraco. Eu não
tenho esse medo. Minha tradição revolucionária diz que aqueles que não
estão dispostos a ir para a cadeia em nome de suas ideias nunca deram
muito valor a elas. Além disso, há lugares piores do que a cadeia. Se eu
tiver de pagar pelas minhas atividades em defesa do controle de
natalidade, uma coisa é certa, o movimento de controle de natalidade não
será interrompido, como tampouco serei impedida de levar adiante a
agitação do controle de natalidade. Se eu me abstive de discutir aqui os
métodos,\footnote{Entre os anos de 1915 e 1916, Goldman deu uma série de
  palestras, por todo os Estados Unidos, para ensinar, sobretudo às
  mulheres da classe trabalhadora, técnicas contraceptivas. Vale lembrar
  aqui da sua formação técnica em enfermagem.} não é porque tenha medo
de uma segunda ordem de prisão, mas porque pela primeira vez na história
da América, o tema do controle de natalidade está circulando, de modo
claro, por meio da informação oral e como eu quero que seja combatido
pelos seus próprios méritos, não desejo dar aqui, às autoridades, a
oportunidade de confundir o que estaria dizendo com alguma outra coisa.
E nesse sentido, quero destacar a completa estupidez da lei. Tenho em
mãos o testemunho dado por policiais que, segundo a declaração deles, é
uma transcrição exata do que falei da plataforma. Esses homens são tão
ignorantes que não transcreveram nenhum método contraceptivo
corretamente. \footnote{Aqui Goldman está se referindo à audiência
  preliminar ao julgamento. Segundo um artigo publicado, em março, no
  \emph{Mother Earth}, intitulado ``Minha prisão e audiência
  preliminar'', Goldman relata que no testemunho dos policiais que
  presenciaram a fala que lhe conduziu à prisão, todos os métodos
  contraceptivos foram descritos erroneamente, embora alegassem
  fidelidade ao seu discurso.} Está perfeitamente dentro da lei que
policiais deem testemunho, mas não está dentro da lei que eu leia o
testemunho que resultou na minha acusação. Vocês podem me culpar por eu
ser anarquista e não respeitar as leis? Também quero pontuar a estupidez
da corte americana. Supostamente, a justiça é para ser encontrada lá.
Supostamente não há procedimentos do tipo da Câmara Estrelada sob a
democracia,\footnote{Câmera estrelada, no original ``star chambre'', é
  uma expressão idiomática utilizada para indicar julgamentos injustos
  resultantes de conspirações e perseguições.} não obstante no dia que
os policiais deram o seu testemunho, isso foi feito entre sussurros ao
juiz como se se tratasse de um confessionário da Igreja Católica e sob
nenhuma circunstância foi permitido às damas então presentes ouvir algo
do que estava ocorrendo.\footnote{Aqui, Goldman está ironizando o fato
  de que quando os policiais testemunharam sobre os métodos
  contraceptivos ensinados na sua palestra, fizeram isso aos sussurros
  para o juiz, de modo que as mulheres presentes na audiência pública
  não pudessem escutá-los.} Que farsa tudo isso! E ainda esperavam que
nós respeitássemos, que obedecêssemos a isso, que nos submetêssemos.

Não sei quantos de vocês estão dispostos a respeitar e obedecer, eu sei
que eu não estou. Permaneço como uma das responsáveis por um movimento
mundial. Um movimento que almeja tornar as mulheres livres do jugo
terrível que é a servidão pela gravidez forçada; um movimento que exige
o direito de que toda criança seja bem-nascida; um movimento que
pretende libertar a mão de obra da sua eterna dependência; um movimento
que almeja dar lugar a um novo tipo de maternidade. Eu considero este
movimento importante e vital o suficiente para que desafie todas as leis
contidas nos livros de estatutos. Eu acredito que este movimento abrirá
caminho não apenas para a discussão livre sobre os métodos
contraceptivos, mas para a liberdade de expressão na Vida, Arte e
Trabalho, e para o direito das ciências médicas manipularem os
contraceptivos como o fazem para o tratamento da tuberculose ou de
qualquer outra doença.

Posso ser presa, posso ser julgada e trancada na cadeia, mas nunca
ficarei em silêncio; eu nunca vou concordar em me submeter à autoridade,
nem farei as pazes com um sistema que degrada a mulher à condição de
mera incubadora, engordando-a com vítimas inocentes. Eu aqui e agora
declaro guerra contra este sistema e não irei descansar até que o
caminho esteja aberto para uma maternidade livre e saudável e uma
infância alegre e feliz.

\chapter[Novamente o movimento do controle de natalidade (1916)]{Novamente o movimento do controle de natalidade\footnote[*]{Texto publicado
  no \emph{Mother Earth}, sete meses após ``Os aspectos sociais do
  controle de natalidade'', em novembro de 1916.}}

Se alguém tiver dúvida sobre a magnitude do crescimento do movimento do
Controle de Natalidade, dois acontecimentos recentes na cidade de Nova
York devem dissipar essa dúvida.

Um deles é a declaração proferida pelo juiz Wadhams da \emph{General
Sessions}, e o outro é o método desesperado que o Departamento de
Polícia de Nova York tem utilizado para lidar com os ativistas do
Controle de Natalidade. A polícia não apenas prende todos que discutem
sobre o assunto abertamente ou distribuem informações sobre o Controle
da Natalidade, mas também plantam acusações contra vítimas inocentes.
Obviamente, o perjúrio não é uma novidade no nosso Departamento de
Polícia; não é de causar surpresa que esse método antiquado esteja sendo
utilizado novamente da maneira mais flagrante possível.

Considere-se, primeiramente, o juiz Wadhams. Uma mulher foi posta diante
dele pelo crime de roubo. A senhora Schnur declarou que foi obrigada a
roubar para conseguir algum pão para os seus seis filhos, dos quais o
mais novo tinha dez meses de idade.

Até cinco anos atrás, Samuel Schnur sustentava sua família com seu
salário de operário num atelier de casacos de crianças do lado leste da
cidade. O confinamento constante acabou afetando o homem: ele pegou
tuberculose. Apesar da doença, continuou trabalhando, até que ela foi
descoberta por um inspetor do Departamento de Saúde, que o impediu de
trabalhar no atelier enquanto estiver com a doença. Desde então, ele não
conseguiu arranjar mais nenhum emprego.

O ônus de sustentar a casa recaiu sobre a sra. Schnur, que passou a
garantir a sobrevivência da família fazendo todo tipo de trabalho.
Recentemente, ela não foi bem-sucedida em conseguir um emprego e gastou
rapidamente o pouco dinheiro que tinha. No mês passado, entrou na casa
de Morris Moskowitz, na rua East Seventh, número 203, e roubou uma
pequena quantia em dinheiro e um relógio, o que resultou na sua prisão.

Segundo as palavras do juiz Wadhams, a sra. Schnur havia sido condenada,
numa ocasião anterior à do roubo, decorrente das mesmas condições, mas,
a sentença foi suspensa. Embora a mulher pudesse ter sido sentenciada a
uma prisão de longo prazo, o tribunal disse que as circunstâncias
incomuns eram tais que justificavam mais um pedido de clemência.

Depois de discutir a condição do marido e sua impossibilidade de cuidar
da família, o juiz Wadhams fez a seguinte declaração:

``Apesar de tudo, ele continua se tornando pai de filhos que, por conta
das condições, têm pouquíssimas chances de serem outra coisa que não
tuberculosos, repetindo, ao crescer, o mesmo processo na sociedade. Não
há lei contra isso.

O problema não é apenas que não tenhamos uma regulação da natalidade
para esses casos, mas também o de que pessoas estão sendo levadas ao
tribunal por fornecer informações para o controle dos nascimentos.
Existe uma lei que elas violam.''

O juiz Wadhams destacou que muitas nações da Europa adotaram a regulação
da natalidade com resultados aparentemente excelentes. Ele perguntou se
os americanos estavam sendo razoáveis em relação ao assunto, como seria
o desejado.

``Creio'', continuou a Opinião, ``que estamos vivendo numa era de
ignorância que, em algum tempo futuro, será vista com o mesmo horror com
que agora lembrarmos do do que antes perseguido e agora nos orgulhamos
que exista. Eis que então, diante de nós, uma família que não para de
crescer em número, embora o marido seja tuberculoso, a mulher traga um
filho no peito e outros filhos pequenos debaixo da saia e não tenham
dinheiro nenhum. ''

Com certeza, valeu a pena ir para a prisão já que, com isso, foi
possível ensinar a um juiz a importância do controle de natalidade para
as massas. No entanto, ir para a cadeia nunca ensina nada à polícia.

Sexta-feira, 20 de outubro, fui intimada a comparecer como testemunha do
Sr. Bolton Hall, no seu julgamento perante as \emph{Special Sessions},
Departamento Seis, por ter distribuído circulares de Controle de
Natalidade num encontro do movimento na \emph{Union Square}, no dia 20
de maio. Hall foi absolvido. Em conjunto com vários amigos, saí do
tribunal por volta das 17 horas e mal cheguei na calçada fui presa pelo
policial Price. Quando lhe foi pedido para mostrar o mandado, ele disse
que era desnecessário, que eu estava sob sua incumbência e que teria de
o acompanhar. Por experiências passadas, já sabia que um policial, caso
dos Centenas Negras da Rússia, é uma espécie de poder absoluto e, assim,
segui com ele até a delegacia da rua Elizabeth. Lá, me foi arbitrada a
fiança de US\$ 1.000 por ter distribuído folhetos de Controle de
Natalidade na mesma reunião na \emph{Union Square} no dia 20 de maio.
Evidentemente, os policiais não levaram em consideração o testemunho
esmagador prestado em nome do Sr. Bolton Hall, um testemunho que, é
claro, também será prestado em meu favor, de que nem o Sr. Hall e nem eu
distribuímos folhetos de Controle de Natalidade. Os policiais decidiram
se envolver numa tramoia e começaram, imediatamente, a colocá-la em
ação.

Vocês devem se recordar que em abril passado, fui presa por \emph{ter
dado informações sobre controle de natalidade}; que eu fui julgada e
considerada culpada e que preferi ir para a cadeia do condado de Queen
ao invés de pagar a fiança de US\$ 100,00. Com isso em vista, não é
necessário enfatizar o quanto acredito na questão do controle da
natalidade e na necessidade de dar às pessoas esse tipo de informação.
Em outras palavras, estou disposta a assumir as consequências se sou
culpada perante o que apetece à lei chamar de ofensa. No entanto, o fato
é que eu não distribuí as circulares em questão e, é óbvio que não
pretendo ser presa e ficar trancafiada na cadeia simplesmente porque os
policiais de Nova York decidiram se coroar com os louros do estancamento
da agitação do Controle de Natalidade.

\asterisc

No dia 27 de outubro, uma sexta-feira, compareci perante o juiz Barlow,
um tipo que lembra os personagens de Dickens ou Victor Hugo. Severo,
pomposo e estúpido. Renunciei a audiência preliminar e serei levada a
julgamento.

No momento, não posso dizer quando o julgamento ocorrerá, mas espero
conseguir um adiamento e possivelmente um julgamento por júri. Enquanto
isso, foi-me arbitrada essa fiança de US\$ 1.000.\footnote{Em janeiro de
  1917, Goldman é absolvida da acusação de ter distribuído panfletos de
  controle de natalidade em 20 de maio de 1916\emph{.}}

Em 30 de outubro, três casos foram julgados no Tribunal Especial
{[}\emph{Court of Special Sessions}{]}: Jesse Ashley e dois jovens da
I.W.W.\footnote{Sigla da organização \emph{Industrial Workers of the
  World} {[}Sindicato dos Trabalhadores Industriais do Mundo{]}.}, Kerr
e Marman. Claro que todos foram considerados culpados. Jesse Ashley foi
condenada a pagar a multa de US\$ 50,00 ou cumprir 10 dias de prisão.
Ela teria preferido a cadeia, mas foi encorajada a abrir um precedente.
Por conta disso, ela pagou a quantia sob protesto e irá recorrer.

Kerr e Marman provavelmente se sairão pior, já que o policial
testemunhou que, enquanto conversavam na Madison Square, os meninos
ofereceram um panfleto de Controle de Natalidade por 10 centavos de
dólar e lhe deram -- Oh, horrores! - ``A preparação militar nos conduz
ao massacre universal'', de E. G. Sequer nos nossos sonhos, poderíamos
imaginar, quando publicamos o panfleto, que papel importante ele iria
desempenhar num tribunal de Nova York.

Os dois meninos insistiram que o que eles fizeram foi vender o panfleto
``A preparação'' e dar os folhetos sobre o Controle de Natalidade de
graça. Muito embora os juízes todos saibam que os policiais nunca
hesitam em cometer perjúrio, Kerr e Marman foram considerados culpados.

É evidente que estamos conseguindo educar os juízes. Não por acaso, um
deles foi muito enfático ao afirmar que pretende distinguir entre
aqueles que dão informações gratuitas sobre Controle de Natalidade, por
convicção, e aqueles que as vendem. Nenhuma consideração dessa natureza
foi feita no caso de Bill Sanger, de Ben Reitman ou mesmo no meu, embora
nenhum de nós tenha vendido informações.

Essa novidade na opinião de um juiz prova, portanto, que a ação direta é
a única ação que conta. Mas para todos nós que desafiamos a lei, o
Controle de Natalidade ainda é uma proposta de salão, apesar de todas as
ligas de Controle de Natalidade neste país.\footnote{Goldman organizou o
  as ligas também com a intenção de combater as leis que proibiam a
  pratica e informação sobre o controle de natalidade. A estragégia era
  a de que assim que um dos membros fosse preso, outro saíriam para
  distruir, de modo a também ser preso e, assim, sucessivamente. A ideia
  era a de que uma série de prisões sem sentido demonstrariam a
  obsolescência da lei.}

\chapter[O camaleônico sufrágio feminino (1917)]{O camaleônico sufrágio feminino\footnote[*]{Texto originalmente publicado
  na revista anarquista ``\emph{Mother Earth}'' (1917).}}

Há quase meio século, as líderes do sufrágio feminino têm apregoado os
resultados milagrosos que se seguiriam à alforria da mulher. Todos os
males sociais e econômicos dos séculos passados seriam abolidos tão logo
a mulher obtivesse o direito ao voto. Todos os erros e injustiças, todos
os crimes e horrores cometidos ao longo das eras seriam eliminados da
vida através do decreto mágico contido num pedaço de papel.

Quando as líderes do movimento eram alertadas para o fato de que tais
alegações extravagantes não convenciam ninguém, elas diziam: ``Esperem
até que tenhamos a oportunidade; esperem até que estejamos frente a
frente com uma grande provação e, então, vocês verão o quão superior é a
mulher em sua atitude para com o progresso social''.

Os oponentes realmente inteligentes do sufrágio feminino -- que estavam
fundamentados na compreensão de que o sistema representativo serviu
apenas para roubar a independência do homem, e que faria o mesmo com a
mulher -- sabiam muito bem que, em lugar nenhum, o sufrágio feminino
exerceu a mínima influência sobre a vida social e econômica das pessoas.
Ainda assim, eles se dispuseram a dar aos expoentes do sufrágio o
benefício da dúvida. Eles estavam dispostos a acreditar que as
sufragistas eram sinceras na sua afirmação de que a mulher nunca seria
culpada das estupidezes e crueldades cometidas pelo homem. Eles
vislumbraram, especialmente nas sufragistas militantes da Inglaterra, um
tipo superior de feminilidade. Não foi a senhora Emmeline
Pankhurst\footnote{Nos anos de 1909 e 1913, a fundadora do \emph{Women's
  Social and Political Union}, ou simplesmente WSPU, Emmeline Pankhurst
  fez uma série de conferências muito bem-sucedidas nos Estados Unidos,
  daí a alusão de Goldman.}, quem fez a afirmação ousada, do alto de uma
plataforma nos Estados Unidos, de que a mulher é mais humana do que o
homem e que nunca seria culpada pelos seus crimes; até porque, para
começar, a mulher não acredita na guerra e nunca apoiará as guerras?

Mas políticos permanecem políticos. Tão logo a Inglaterra entrou na
guerra, por razões humanitárias, é claro, as damas sufragistas
esqueceram imediatamente toda a bazófia da superioridade e bondade da
mulher e imolaram o seu próprio partido no altar do mesmo governo que
rasgou suas roupas, puxou seus cabelos e as alimentou à força\footnote{A
  greve de fome era uma das formas de protesto mais utilizadas pelas
  sufragistas da WSPU -- especialmente, quando presas pelos atos em nome
  da causa que, em geral, eram criminalizados como vandalismo. Muitas
  sufragistas de renome escreveram relatos sobre a experiência de serem
  alimentadas à força nas prisões, denunciando os abusos físicos e
  emocionais que acompanhavam a prática. O uso repetido da sonda
  nasogástrica para a alimentação forçada causava, por si só, danos
  permanentes à saúde. Em alguns relatos, a analogia entre a alimentação
  forçada e o estupro, embora não exposta diretamente, é suficientemente
  clara. Em 1913, foi promulgada a lei que ficou conhecida como Lei do
  Gato e Rato {[}\emph{Cat and Mouse Act}{]}, em que já não mais se
  imporia a alimentação forçada às grevistas, mas ao invés disso, as
  deixaria chegar ao limite das suas forças no presídio, até que sob o
  risco de morte fossem temporariamente libertadas para a sua
  recuperação -- o que, diga-se de passagem, isentava o governo de
  qualquer responsabilidade para com eventuais mortes decorrentes das
  greves prolongadas. A imagem da sufragista, voluntariamente em greve
  de fome, em sua cela isolada, teve forte ressonância na imaginação
  pública.} por conta das suas atividades militantes. A senhora
Pankhurst e o seu bando tornaram-se mais apaixonadas na sua mania de
guerra e na sua sede pelo sangue do inimigo do que os militaristas mais
insensíveis.\footnote{Com a entrada da Inglaterra na Primeira Guerra
  Mundial, Pankhurst e a sua filha mais velha Christabel, na época a
  líder formal da WSPU, comunicaram, em agosto de 1914, a decisão de
  interromper todas as atividades de militância do grupo até o final da
  guerra -- o que, como não poderia ser diferente, conduziu ao
  desligamento e formação de novas células por parte de uma minoria que
  não estava de acordo. Como resposta a este apoio inesperado,
  incialmente visto com desconfiança, o governo britânico libertou todas
  as sufragistas presas e retirou todas as queixas, o que possibilitou a
  ambas as Pankhursts, mãe e filha, retornarem à Inglaterra sem correr o
  risco de serem, mais uma vez, presas. Emmeline e Christabel
  transformaram a WSPU em veículo de apoio à guerra e passaram a se
  apresentar como patriotas feministas, defendendo a importância da
  participação das mulheres na guerra em defesa da esclarecida
  democracia britânica, então ameaçada pela autocracia militarista da
  Alemanha. Naquele momento, segundo elas, o dever das feministas havia
  se convertido em inspirar o patriotismo.} Elas consagraram tudo, até a
própria sensualidade, como meio de atrair os homens ainda relutantes
para o campo militar, para as trincheiras e a morte.\footnote{Dado que,
  na época, o recrutamento militar na Inglaterra era voluntário e não
  obrigatório, Emmeline passou a incentivar nos seus comícios, após
  agosto de 1914, o alistamento dos homens na guerra, tornando-se uma
  espécie de agente de recrutamento não oficial. Apesar da defesa da
  participação das mulheres na guerra, reforçou estereótipos ao ecoar a
  visão comum acerca da inadequação das mulheres à condição de soldado,
  posto ser o mais alto dever destas o trabalho de manutenção da
  sociedade e sobretudo a maternidade, o que seria incompatível com uma
  presença efetiva nas trincheiras. Às mulheres caberia impedir o
  colapso da nação, ao assumir os empregos tradicionalmente ocupados por
  homens que, assim, estariam livres para lutar em defesa do seu país e
  das democracias da Europa e, portanto, em defesa das suas mulheres --
  conforme sugeria marotamente nas suas conferências, um dos pontos que
  parece estar sendo aqui denunciado por Goldman. Ainda que a defesa da
  assunção das mulheres de trabalhos até então exclusivamente destinados
  aos homens fosse de encontro do lugar tradicionalmente reservado à
  mulher na sociedade britânica, já não se tratava mais de subversão. Em
  realidade, no final da guerra, as Pankhursts, mãe e filha mais velha,
  passaram da condição de ``terroristas'' a ultrapatriotas defensoras do
  \emph{status quo.} Para a surpresa de muitos, em 1926, Emmeline veio a
  se filiar ao Partido Conservador do Reino Unido, correndo como
  candidata do partido, em 1928 cuja campanha teve de ser abortada por
  conta de um escândalo envolvendo outra das suas filhas, dissidente do
  WSPU -- que havia tido um filho fora do casamento e dado declarações
  públicas em favor da sua condição de mulher e mãe ``livre''.} Por tudo
isso, agora elas serão recompensadas ​​com o sufrágio. Até Asquith, o
antigo inimigo da organização de Pankhurst, está agora convencido de que
a mulher deve ter direito ao voto, já que se provou tão feroz no seu
ódio e tão empenhada na conquista.\footnote{Herbert Henry Asquith ficou
  conhecido pela oposição atuante e agressiva contra o sufrágio feminino
  -- iniciada em 1892 e persistindo ao longo de boa parte do seu mandato
  como primeiro-ministro do Reino Unido (1908-1916). Embora o seu
  posicionamento contra o sufrágio feminino tenha dado indícios de
  mudança já em 1912, em 1915 declarou publicamente que dado o
  engajamento demonstrado, já não poderia negar a reivindicação das
  mulheres uma vez finda a Guerra -- declaração que, ao que parece, é a
  que Goldman está se referindo aqui. Seja como for, a desconfiança e
  repúdio das duas Pankhursts contra Asquith, tornadas públicas sempre
  que tinham oportunidade, não cessou até a sua renúncia em 1916. O
  sufrágio feminino no Reino Unido foi garantido por lei apenas em 1918,
  não obstante apenas para uma pequena parcela de mulheres, cujo projeto
  já havia sido denunciado por Goldman no texto ``O sufrágio feminino''
  de 1910.} Todos saúdam as inglesas que compraram seu direito ao voto
com o sangue de milhões de homens já sacrificados ao monstro Guerra. O
preço é de fato alto, mas também altos serão os cargos reservados para
as damas da política.

O partido sufragista americano, desprovido de qualquer ideia original
desde os dias de Elizabeth Cady Stanton, Lucy Stone e Susan
Anthony,\footnote{Elizabeth Cady Stanton (1815-1902), Lucy Stone
  (1818-1893) e Susan Anthony (1820-1906) têm em comum não apenas o fato
  de terem sido pioneiras do movimento sufragista estadunidense, ficando
  conhecidas como o ``Triunvirato'' do movimento sufragista, quanto
  também de serem militantes aguerridas do abolicionismo. As três não só
  se conheceram como trabalharam juntas, influenciando-se mutuamente. Há
  quem defenda que as primeiras sufragistas tinham ligações diretas com
  alguns dos povos nativos norte-americanos, de modo que as suas ideias
  feministas refletiam, em muitos aspectos, os papéis de liderança e
  poder desempenhados pela mulher nesses povos. Em 1866, o
  ``Triunvirato'' fundou a \emph{American Equal Rights Association}, que
  militava a favor da igualdade de direitos e sufrágio para todos os
  estadunidenses independentemente da raça, cor ou sexo. A associação,
  porém, não teve vida longa: em 1869, Anthony e Stanton fundaram a
  \emph{National Woman Suffrage Association} (NWSA) e, poucos meses
  depois, Stone veio a fundar a \emph{American} \emph{Woman Suffrage
  Association} (AWSA). A gota d'água que culminou na cisão foi a
  discussão em torno do apoio ou não das sufragistas à Décima Quinta
  Emenda que garantia a todos os cidadãos estadunidenses \emph{homens},
  independentemente da cor e raça, o direito ao voto. O motivo do
  desacordo estava relacionado às táticas de ação e não aos fins
  propriamente ditos. Em 1890, os dois grupos, então rivais em muitas
  pautas, fundiram-se, após longa negociação, na organização
  \emph{National American} \emph{Woman Suffrage Association} (NAWSA).
  Antes da fusão, em 1889, líderes de ambas as associações, o que
  incluía o ``Triunvirato'', publicaram juntas uma ``Carta aberta para
  as mulheres da América''.} irá necessariamente macaquear, repetir, de
modo estúpido, o exemplo dado por suas irmãs inglesas. Nos dias heroicos
da sua militância, a sra. Pankhurst e suas seguidoras foram
completamente repudiadas pelo partido sufragista americano. Uma dama
elegante e respeitável como a senhora Catt\footnote{Carrie Chapman
  Catt~sucedeu Susan Anthony na presidência da \emph{National American}
  \emph{Woman Suffrage Association}, em 1900. Ela exerceu dois mandatos
  como presidente da NAWSA, o primeiro de 1900 a 1904 e o segundo de
  1915 a 1920, sendo, portanto, presidente no ano em que o presente
  artigo foi redigido e publicado. Sob a liderança de Catt, a NAWSA
  pautava a sua defesa do sufrágio feminino na concepção ampla e
  repetidamente combatida por Goldman: a de que as mulheres,
  essencialmente diferentes dos homens, trariam as suas virtudes
  domésticas de modo a purificar a política. Muitas décadas após a sua
  morte, nos 1990, Carrie Chapman Catt teve o seu nome associado ao
  racismo; foram-lhe atribuídas uma série de declarações racistas e
  xenófobas, dentre elas a de que o sufrágio feminino fortaleceria a
  causa dos supremacistas brancos. A discussão em torno do racismo de
  Catt culminou num movimento liderado por estudantes da \emph{Iowa
  State University,} ainda na década de 1990, com o objetivo de remover
  o seu nome de um dos edifícios da instituição -- o que não ocorreu. A
  factualidade de boa parte das declarações foi contestada por
  pesquisadores que alegaram ausência de registros históricos que
  confirmassem as fontes secundárias. No que diz respeito à infame
  associação entre o sufrágio feminino e a supremacia branca, deveras
  contida num ensaio de Catt, também datado de 1917 e intitulado: ``O
  sufrágio feminino na emenda federal'' {[}\emph{Woman Suffrage by
  Federal Amendment}{]}, seus estudiosos chamam a atenção para a
  necessidade de interpretação do contexto em que tal associação é
  estabelecida -- dentro e fora do texto. Até hoje, a discussão perdura.}
não poderia mesmo ter algo a ver com facínoras como são todos os
militantes. No entanto, quando as sufragistas da Inglaterra, de olho nos
privilégios luxuriosos do Parlamento, deram suas cambalhotas e saltos
mortais, o partido sufragista americano seguiu o exemplo. A verdade é
que a sra. Catt sequer esperou que a guerra fosse declarada pelo seu
país. Ela se saiu ainda melhor do que a Sra. Pankhurst: ofereceu o seu
partido ao militarismo, prometeu apoio a todas as medidas autocráticas
do governo muito antes que houvesse necessidade de tudo
aquilo.\footnote{Ainda antes da entrada formal dos Estados Unidos na
  guerra, Catt, apesar de pacifista e membro de organizações pacifistas,
  rapidamente, seguindo o exemplo dado pelas Pankhursts-- segundo a tese
  de Goldman aqui exposta --, resolveu colocar a força da NAWSA à
  serviço da guerra, fosse para a mobilização de apoio popular; para a
  arrecadação de fundos e alimentos; preparação e suporte às mulheres
  que viessem a ocupar os empregos deixados pelos homens etc. Muito
  embora o presidente Wilson se opusesse ao sufrágio feminino desde o
  início do seu mandato, em 1913, a editora da NWASA reimprimiu as
  mensagens de guerra do presidente, valendo-se delas nos seus esforços
  de mobilização da população. O mote legitimador era a de que aqueles
  que já estão treinados em lutar pela democracia no seu país, estão
  prontos para lutar pela democracia mundial. Como as Pankhursts, Catt,
  e muitas outras sufragistas norte-americanas, passaram a vincular a
  causa sufragista ao patriotismo e às virtudes cívicas.} Por que não?
Por que desperdiçar mais cinquenta anos fazendo lobby pelo direito ao
voto se é possível obtê-lo por meio da simples traição a um ideal? Que
ideais podem subsistir entre políticos, afinal!

O argumento dos opositores, de que a mulher não precisa do voto porque
ela tem uma arma mais forte -- o seu sexo --, havia sido respondido com
a declaração de que o voto justamente libertaria a mulher da necessidade
degradante do apelo sexual. Como é possível que esse orgulho ostensivo
tenha dado origem à campanha, agora iniciada pelo partido sufragista, de
atrair a masculinidade da América para o mar de sangue europeu? As
justíssimas integrantes do partido sufragista não estão apenas abordando
descaradamente cada jovem e cada homem, persuadindo-os a se alistar,
como estão induzindo as esposas e namoradas a jogar com as emoções e
sentimentos dos seus companheiros, a fim de incitar o seu sacrifício ao
Moloch do Patriotismo e da Guerra.\footnote{Através de uma ampla
  campanha encabeçada pela NAWSA, o governo de Wilson inaugurou, em
  1917, o \emph{Woman's Committee of Council Nacional Defense}, que
  subordinado ao \emph{Council Nacional Defense}, por sua vez
  subordinado Departamento de Guerra, teria a responsabilidade de
  organizar toda a mobilização civil durante a guerra.}

Como essa tarefa poderá ser bem-sucedida? Certamente, não pela via do
argumento. Se durante os últimos cinquenta anos as mulheres militantes
falharam em convencer, a maioria dos homens, de que a mulher tem direito
à igualdade política, elas certamente não os convencerão,
repentinamente, de que devem ir para a morte certa, enquanto, elas,
mulheres permanecerão, em segurança, enfiadas nas suas casas, a costurar
ataduras. Não, não foi nenhum argumento, razão ou humanitarismo que o
partido sufragista prometeu ao governo; e sim, o poder da atração
sexual, o apelo vulgar, persuasivo e envolvente da mulher liberada a
serviço da glória do seu país. Que homem pode resistir a isso? Se muitos
dos grandes homens perderam a sua sanidade e capacidade de julgamento,
quando entorpecidos pelo apelo sexual, como a juventude da América
poderia resistir?

O rei está nu. As damas sufragistas já deram provas de que a sua
prerrogativa não é nem a inteligência, nem a sinceridade, e que a sua
ostentação de igualdade está completamente podre; que, inclusive, na
luta pelo direito ao voto, o apelo sexual é o seu único recurso, e a
recompensa da politicagem mais chula, seu único objetivo. Elas, agora
estão utilizando ambos, recurso e objetivo, para alimentar o cruel
monstro da guerra, embora devam saber que, por mais terrível que seja o
preço pago pelo homem, não é nada quando comparado à crueldade,
brutalidade e ultraje aos quais a mulher está sujeita durante a guerra.

O crime que as líderes do partido do sufrágio feminino americano estão
cometendo contra o seu eleitorado é análogo à relação do alcoviteiro com
sua vítima. A maioria dessas lideranças já está muito velha para exercer
qualquer efeito sobre o alistamento via o apelo sexual ou mesmo para
prestar qualquer serviço pessoal ao seu país. Ao prometerem ao governo o
apoio do partido, estão vitimando os membros mais jovens. Isso pode soar
excessivamente cruel, mas, no entanto, é a verdade. De outro modo, como
poderíamos explicar o compromisso de bater de porta em porta para
trabalhar a histeria patriótica nas mulheres, que, em troca, devem
utilizar a sua sexualidade para fazer seus homens se alistarem? Em
outras palavras, o mesmo atributo que a mulher foi forçada a usar para
garantir o seu \emph{status} econômico e social na sociedade, e que as
damas sufragistas sempre repudiaram, está agora começando a ser
explorado a serviço do Senhor da Guerra.\footnote{Não foram poucas as
  líderes sufragistas que, nos seus esforços no tempo da guerra, foram
  tomadas de grande fervor patriótico.}

Para fazer justiça ao \emph{Woman's Political Congressional Union} e a
alguns poucos membros do partido sufragista, é preciso dizer que há
resistência à persuasão das líderes sufragistas. Infelizmente, porém,
para dizer a verdade, o \emph{Woman's Political Congressional Union}
está em cima do muro: não se posiciona nem a favor da guerra, nem da
paz.\footnote{Aqui Goldman está se referindo ao \emph{National Women's
  Party}, conforme foi reorganizado, em 1916, o \emph{Woman's Political
  Congressional Union}. Diferente da NAWSA, de onde muitas das suas
  lideranças haviam sido expulsas, as mulheres do \emph{National Women's
  Party} continuaram a protestar contra a oposição do governo de Wilson
  ao sufrágio feminino, independentemente da guerra. Em realidade,
  tomaram como argumento central a contradição de um governo que estaria
  a lutar para assegurar a democracia no exterior, enquanto negava às
  mulheres cidadãs do seu país o direito ao voto. Algumas foram presas
  nesses protestos. Ao se concentrarem, de modo praticamente exclusivo,
  na causa do sufrágio feminino, recusaram qualquer posicionamento
  acerca da atuação dos Estados Unidos na guerra.} Não haveria problema
nisso, desde que o monstro se deslocasse exclusivamente por sobre a
Europa. Agora que está se alastrando em casa, o \emph{Congressional
Union} descobrirá que silêncio é sinal de consentimento. A recusa em se
posicionar de modo taxativo contra a guerra praticamente torna os seus
membros uma parte dela.

Em meio a essa confusão entre as facções do movimento sufragista, é
realmente refrescante encontrar uma mulher decidida e firme. A recusa de
Jannette Rankin em apoiar a guerra faz mais, no sentido de aproximar a
mulher da emancipação, do que todas as medidas políticas juntas. No
momento presente, ela é sem dúvida considerada um anátema, uma traidora
do seu país.\footnote{Primeira mulher eleita para a Câmera dos Deputados
  nos Estados Unidos, em 1916, Jannette Rankin, até então um dos
  expoentes da NAWSA, pagou um preço alto ao escolher manter a coerência
  entre o pacifismo pelo qual havia militado ao longo de toda vida, e a
  sua atitude política: votou contra a declaração de guerra à Alemanha
  em 1917.} Mas isso não deveria entristecer a senhorita Rankin. Todos
os homens e mulheres valiosos foram denunciados como tal. No entanto,
eles e não os patriotas gritalhões e covardes são efetivamente valorosos
para a posteridade.

\chapter[Louise Michel, uma refutação (1923)]{Louise Michel, uma refutação\footnote[*]{Carta aberta dirigida ao médico e
  se Magnus Hischfeld, escrita em inglês e publicada em alemão na
  \emph{Jahrbuch für sexualle Zwischenstufen}, como resposta ao ensaio
  de Karl von Levetzow, no qual ele defende que a anarquista francesa
  Louise Michel seria lésbica. O artigo de von Levetzow foi publicado na
  mesma revista, no ano da morte de Michel, em 1905. Uma tradução para o
  inglês da carta foi publicada, em 1976, na coletânea de textos e
  documentos organizada por Jonathan Katz, intitulada \emph{Gay American
  History}. Aqui, estamos traduzindo do rascunho da carta original em
  inglês (1923).}}



Prezado Dr. Hirschfeld:

Eu tenho conhecimento do seu excelente trabalho sobre psicologia sexual
já há alguns anos. Sou admiradora da sua luta corajosa pelos direitos
das pessoas que, pela sua própria natureza, não se expressam sexualmente
do modo comumente considerado ``normal''.\footnote{A \emph{Jahrbuch für
  sexualle Zwischenstufen} era a revista oficial do \emph{Scientic
  Humanitarian Committee}, a primeira organização dedicada a promover os
  direitos dos homossexuais e transgêneros. Dr. Hirschfeld, era não só
  editor da revista, como fundador do \emph{Scientic Humanitarian
  Committee}.} E agora que tive a sorte de conhecê-lo e acompanhar os
seus esforços de perto, estou mais do que nunca impressionada com a sua
personalidade e com o espírito que tem lhe sustentado nessa tarefa tão
difícil. Sua prontidão em publicar a minha refutação à avaliação de
Frhr. von Levetzow de que Louise Michel seria uraniana\footnote{Aqui,
  Goldman utiliza uma expressão então corrente na época, no original
  ``\emph{urning}''. Em 1889, no volume de lançamento do \emph{Jahrbuch
  für sexualle Zwischenstufen}, Hirschfeld publicou uma série de cartas
  do alemão Karl Heirich Ulrichs, escritor de romances homossexuais, e
  considerado o primeiro ativista a lutar politicamente pelos direitos
  dos homossexuais. Embora fosse advogado, por conta das leis
  discriminatórias contra os homossexuais, Ulrich morreu em miséria, na
  Itália. Ele foi responsável por cunhar diversos termos para designar
  diferentes formas de sexualidade, dentre eles, o termo ``uraniano''
  que indica uma ``alma feminina no corpo de um homem''. Conforme se
  verá, Goldman intercambia livremente as palavras ``uraniano'' e
  ``homossexual''.} prova, se provas são necessárias, que você tem um
senso de justiça elevado cuja meta exclusiva é a averiguação da verdade.
Agradeço por isso e pela postura habilidosa e heroica que você tem
adotado contra a ignorância e a hipocrisia, a favor do esclarecimento e
do humanismo.

Antes de tratar do artigo, propriamente dito, permita-me dizer o
seguinte: não foi qualquer preconceito contra a homossexualidade ou
aversão aos homossexuais, o que me impeliu a destacar os erros contidos
na argumentação do autor. Se Louise Michel tivesse demonstrado, em algum
momento, traços homossexuais para aqueles que a conheciam e amavam, eu
seria a última pessoa negar esse ``estigma''. Na verdade, considero tal
estigma uma tragédia para aqueles que são sexualmente diferenciados num
mundo tão desprovido de compreensão para com os homossexuais e tão
ignorante no que diz respeito ao significado e à importância de toda a
variedade que envolve o sexo. É certo que não penso que essas pessoas
sejam inferiores, menos morais ou menos capazes de sentimentos e ações.
E acima de tudo: não penso que seja necessário ``limpar'' a minha
ilustre professora e camarada Louise Michel da acusação de
homossexualidade. Seu valor para a humanidade, sua contribuição para a
emancipação dos escravos são tão grandes que, qualquer que fosse o seu
modo de gratificação sexual, em nada poderia lhe acrescentar ou
depreciar.\footnote{A anarquista francesa Louise Michel (1830-1905), ou
  a "Virgem Vermelha", como ficou conhecida, foi a mulher de maior
  destaque na Comuna de Paris. Desempenhou diversos papéis na Comuna,
  inclusive o de soldado, participando ativamente da legendária luta nas
  barricadas. Fosse na prisão, exílio ou quando em liberdade, Louise
  Michel se manteve absolutamente dedicada às atividades revolucionárias
  até a sua morte, em 1905.}

Anos atrás, quando eu não sabia nada sobre psicologia sexual, e o meu
único contato com homossexuais havia sido algumas mulheres que conheci
na prisão, onde estava encarcerada por minhas opiniões políticas, saí em
defesa de Oscar Wilde. Como anarquista, meu lugar sempre foi ao lado dos
perseguidos. E vi na perseguição a Oscar Wilde e na acusação contra ele,
o reflexo da injustiça cruel e da hipocrisia da mesma sociedade que o
enviou ao seu martírio. Foi por isso que o defendi.\footnote{Oscar Wilde
  foi condenado a dois anos de prisão, por sodomia, em 1895.}

Posteriormente, quando fui à Europa, e lá me deparei com as obras de
Havelock Ellis, Krafft-Ebbing, Carpenter e muitos outros, o crime contra
Oscar Wilde e demais homossexuais se apresentou sob uma luz ainda mais
gritante. Desse momento em diante, passei a usar a minha caneta e a
minha voz na defesa daqueles a quem a natureza mesma destinou ser
diferente em sua psicologia e necessidades sexuais. Os seus trabalhos,
prezado doutor, ajudaram-me muitíssimo a iluminar essa questão
extremamente complexa, que é a da psicologia sexual e a humanizar a
atitude das pessoas que passaram a vir me escutar.

Dito isso, os seus leitores poderão ver que não tenho quaisquer
preconceitos, ou a menor antipatia, pelos homossexuais. Muito pelo
contrário. Tenho entre os meus amigos, homens e mulheres, tanto
uranianos completos, quanto bissexuais. E o que encontro neles é uma
inteligência, habilidade, sensibilidade e charme, em geral, muito acima
da média. Sinto-me profundamente ligada a eles, porque sei que o seu
sofrimento é maior e mais complexo do que o da maioria das pessoas. No
entanto, há uma tendência predominante entre os homossexuais e à qual eu
me oponho. É a tentativa de reivindicar toda personalidade excepcional
para o seu credo, de atribuir a essas personalidades traços e
características que são a eles inerentes.

Pode ser uma espécie de condicionamento psicológico que as pessoas que
são perseguidas busquem apoio nos tipos excepcionais das mais diferentes
épocas. A miséria sempre busca companhia. Por exemplo, os judeus
consideram que a maioria dos grandes homens e mulheres do mundo são de
origem judaica ou possuem características judaicas. Os irlandeses fazem
o mesmo. Os hindus lhe dirão que a civilização deles é a melhor do mundo
e assim por diante. O mesmo ocorre com os párias políticos. Os
socialistas reivindicam Walt Whitman e Oscar Wilde para as teorias de
Marx, enquanto muitos anarquistas situam Nietzsche, Wagner, Ibsen e
outros entre os seus. Certamente, a grandiosidade sempre acompanha a
versatilidade, mas considero uma imposição reivindicar alguma pessoa
criativa de destaque para minhas ideias, a menos que ele ou ela as tenha
declarado por si.

Caso se tome esse tipo de alegação característica a muitos homossexuais
como certa, chegar-se-á à conclusão de que não há e nem pode haver
talento e grandiosidade em pessoas que não sejam sexualmente invertidas.
A perseguição gera o sectarismo; o que, em contrapartida, torna as
pessoas limitadas ao seu intuito e, muitas vezes, injustas na avaliação
de outras pessoas. Prefiro pensar que Frhr. Karl von Levetzow sofreu de
uma overdose de sectarismo homossexual. Acrescente-se a isso, a sua
visão bastante antiquada da mulher. Ele vê a mulher apenas como sedutora
de homens e incubadora de crianças e, num sentido mais vulgar, como
cozinheira e faxineira doméstica. Qualquer mulher que não apresente
esses atributos antiquados de feminilidade, será imediatamente
classificada pelo autor como uraniana. À luz das realizações da mulher
moderna nos mais diversos domínios do pensamento humano e da ação
social, esse tipo de visão da mulher, condizente ao homem convencional,
malmente merece ser considerada. Se me propus tratar aqui dessa
compreensão ultrapassada do autor de ``Louise Michel'', foi unicamente
para demonstrar a que conclusões absurdas se pode chegar, quando se
parte de uma premissa absurda.

Minhas críticas a von Levetzow não me impedem de lhe prestar homenagem
como grande artista literário, capaz de compreender, empaticamente, uma
grande alma. Na verdade, sinto-me inclusive culpada por ter de dissecar
o seu artigo. É como se eu tentasse fatiar um grandioso e radiante
retrato pintado à mão por um mestre, pois a imagem de Louise Michel,
traçada pela pena de von Levetzow, é uma obra-prima. Por isso, todos nós
que conhecemos e amamos essa mulher maravilhosa estamos em dívida com
ele. Seja como for, a verdade exige que eu coloque os meus sentimentos
de lado e lide com os fatos.

Para que eu pudesse tratar adequadamente dos pontos levantados no
artigo, seria necessário republicar o texto na íntegra, juntamente com a
minha réplica, ou, no mínimo, fazer uma série de citações extensas,
para, então, abordar cada ponto detalhadamente. O problema é que um tal
feito ocuparia muito espaço do valioso ``Jahrbuch''. Desse modo,
contentar-me-ei com uma simples síntese dos pontos mais relevantes que
foram levantados como prova da homossexualidade de Louise Michel.

Quais são esses pontos?

Primeiro, Louise Michel foi uma criança excepcional: ávida por
conhecimento e problemas científicos, era uma leitora voraz. Segundo, os
seus brinquedos, diferentemente dos das outras meninas, não eram
bonecas, mas sapos, besouros, camundongos e outros seres vivos.
Terceiro, Louise Michel brincava muito com o seu primo (o que, a
propósito, deveria provar que Louise era uma garota perfeitamente
normal, caso contrário escolheria sobretudo garotas para a sua
companhia), escalava árvores, organizava expedições de caça, corria
livremente, enfim estava completamente imersa nas brincadeiras e
travessuras dos meninos. Quarto, ela se tornou totalmente indiferente à
sua aparência, odiava e se opunha aos babados das vestimentas femininas,
espartilhos, sapatos de salto alto e todo o resto; era terrivelmente
displicente consigo mesma e desorganizada em seus hábitos e para com
tudo aquilo que a cercava. Quinto, Louise Michel era extraordinariamente
corajosa, desprovida do sentimento de medo, ousava ao ponto da
imprudência. Seu poder de resistência ao sofrimento físico dificilmente
poderia ser igualado pelos homens. Sexto, com a exceção dos seus
camaradas, nenhum homem fez parte da sua vida. Por outro lado, ela
estava sempre cercada ``de amigas apaixonadas''. Por último, embora não
menos importante, Louise Michel era matemática e compositora, amava
esculturas, era uma entusiasta da música de Wagner e fez muitas coisas
que as mulheres, em geral, nunca fazem. O autor enfatiza bastante o
corpo retangular de Louise Michel, os seus seios pequenos e outras
supostas características masculinas. Em suma, ele utiliza todos os
argumentos imagináveis ​​para provar a masculinidade de Louise Michel,
argumentos que são usados, ​​desde os tempos imemoriais, contra a mulher
sempre que ela busca se elevar da condição de prisioneira do harém e
alcançar a posição de igualdade na vida com o homem.

Vejamos quão verdadeiros são esses chamados ``fatos''.

Primeiro: a inclinação precoce de Louise Michel para os problemas
verdadeiramente profundos da vida, a leitura ávida de livros sérios e o
seu senso matemático são característicos a um número significativo de
mulheres excepcionais; para mencionar apenas algumas, esse é o caso de
Sofia Kovalevskaya, Marie Bashkirtseff e, na atualidade, Madame Curie.
Kovalevskaya resolveu complexos problemas matemáticos aos oito anos de
idade e, quando tinha apenas vinte e cinco anos, já estava entre os
maiores matemáticos da época. Bashkirtseff possuía uma compreensão
psicológica do que se passava em seu entorno muito mais profunda do que
boa parte dos homens mais célebres; ocupou-se do estudo da ciência,
sociologia, literatura, arte, música, aos dez anos de idade, e se tornou
uma das figuras mais marcantes do seu tempo. Madame Curie, por sua vez,
é bastante conhecida, e não como mera auxiliar do marido, mas como uma
autoridade científica independente... No entanto, essas mulheres, e
muitas outras além delas, não só não eram homossexuais, como eram
extremamente femininas; uma feminilidade que foi, em grande medida, o
motivo da maior tragédia de suas vidas, posto que os homens que
conheceram não conseguiram compreender o espírito dessas mulheres,
sequioso do amor e do companheirismo de um homem. Foi assim que Sofia
Kovalevskaya desperdiçou a sua substância com o seu marido e, num
segundo momento, com uma paixão violenta por um compatriota que nunca
pôde suspeitar da chama que estava a consumir essa grande mulher do
século XIX. No que diz respeito a Madame Curie, não desejo entrar na sua
vida privada -- provavelmente, ela consideraria uma invasão. Seja como
for, pelo tanto quanto se sabe sobre a sua vida privada, ela parece
extremamente feminina, sem nenhuma tendência homossexual.

Tenho certeza de que Frhr. von Levetzow nunca presenciou a brincadeira
de garotas americanas, saudáveis e normais. Ele iria perceber que é
possível brincar, escalar árvores, divertir-se com sapos, besouros e
cobras e fazer todo tipo de coisas relacionadas a ``meninos'', e, ainda
assim, se tornar uma mulher extremamente frívola, tipicamente feminina
e, não raro, absolutamente inútil. Por outro lado, existem inúmeras
mulheres americanas de valor imensurável, em quase todas as esferas da
vida, que apesar de muito traquinas na infância, são amantes apaixonadas
dos seus homens, mães dos seus filhos e, ao mesmo tempo, uma grande
força moral nos mais diferentes movimentos que lutam por valores sociais
mais profundos e elevados para o seu país.

Quanto à coragem de Louise Michel, à sua ousadia que beirava a
imprudência, ausência de medo e resistência física impressionante -- é
impossível que eu não concorde que ela, de fato, possuía tais
qualidades. Não obstante, seria injusto para com o sem-número de
revolucionárias russas, se eu enfatizasse todos esses traços
maravilhosos de Louise sem lhes dar o devido crédito. É evidente que o
autor do artigo nunca teve notícias dessas mulheres; para mencionar
apenas algumas: Perovskaya, Gelfman, Figner, Breskovskaya, Kovalskaya,
Volkenstein e, atualmente, Spiridonova.\footnote{Sophia Perovskaya,
  Gesya Gelfman, Vera Figner, Catherine Breshkovskya, Yelizaveta
  Kovalskaya, Ludmilla Volkenstein, and Maria Spiridonova.} Todas essas
mulheres foram heroicas na grande luta revolucionária da Rússia. Elas
realizaram os feitos mais ousados ​​e seguiram para a morte ou para a
Sibéria e a katorga,\footnote{Campos de trabalhos forçados na Sibéria.}
um calvário ainda maior, com um sorriso nos lábios. Perovskaya preferiu
morrer na forca com o seu marido amado do que fugir e garantir a sua
vida, o que ela, facilmente, poderia ter feito. Gelfman e Figner eram
tão profundamente femininas que sofreram mais com a falta de beleza e
delicadeza indispensáveis à vida de uma mulher sensível do que com todos
os horrores físicos da prisão. Kovalevskaya continuou sua luta rebelde
ao longo de todos os anos que passou na prisão -- algo em torno de vinte
e dois anos. Volkenstein e Figner eram extremamente belas e femininas na
sua vida amorosa e nas relações que estabeleciam. No que diz respeito a
Spiridonova, ela foi submetida às torturas mais terríveis, o que inclui
ser estuprada por oficiais russos embriagados e ter o seu corpo nu
queimado com charutos acesos -- muito embora, nem assim, qualquer som
pôde ser ouvido dela. Spiridonova, porém, é uma pessoinha delicada e
frágil, profundamente apaixonada pelos seus camaradas e tão sensível
quanto uma flor. Esses poucos exemplos devem ser suficientes para
convencer qualquer pessoa que não esteja submersa no sectarismo ou nas
velhas noções ultrapassadas da natureza da mulher, de que é possível ser
uma mulher absolutamente feminina e, ao mesmo tempo, uma grande rebelde
e combatente. Nesse sentido, eu poderia seguir enumerando mulheres de
todos os países, de todas as idades e de todos os climas, que se puseram
lado a lado com os homens nas grandes lutas pelos direitos humanos e
pela própria emancipação; mulheres que, certamente, eram tão corajosas e
ousadas quanto os seus companheiros e, ainda assim, não tinham nenhum
traço de masculinidade ou homossexualidade em si.

E isso me traz à conclusão absurda que von Levetzow tirou da
grandiosidade trágica presente no último encontro entre Dombrovski e
Louise Michel nas barricadas. \footnote{Jaroslaw Dombrowski (1836-1871)
  era polonês, e militar do Exército Imperial Russo. Após participar, em
  1863, de um levante mal sucedido pela independência da Polônia, fugiu
  para Paris e lá, estabeleceu conexões com os radicais. Participou
  ativamente da Revolução de 18 de março de 1871, vindo a se tornar o
  comandante militar da Comuna de Paris, pouco antes da invasão do
  exército francês que pôs fim à Comuna. O massacre ficou conhecido como
  ``Semana sangrenta''. Dombrowski morreu de complicações decorrentes
  dos ferimentos sofridos nas barricadas.} O autor é tão limitado pela
sua concepção masculina da mulher que simplesmente não é capaz de
entender como duas pessoas, ante a iminência do colapso da causa que
amavam mais do que a própria vida, poderiam se tratar como camaradas.
Ele observa: ``se Dombrovski tivesse visto uma mulher em Louise, ele
teria lhe dado uma tapinha na bochecha; mas, ao invés disto, ele
estendeu ambas as mãos e apertou as dela no seu último adeus''. Fico
realmente surpresa que um homem com a sensibilidade de von Levetzow seja
capaz de tamanha vulgaridade. Prefiro pensar que isso seja fruto da
ignorância sobre o tipo de relação maravilhosa que existia e ainda
existe entre homens e mulheres envolvidos na luta por um ideal ou que
compartilham uma causa comum. Pois a verdade é que, muitas vezes, a
consciência da diferença de sexo é simplesmente obliterada entre eles;
eles são camaradas, capazes do mais alto sacrifício um pelo outro e de
grande devoção um ao outro. Aqui, mais uma vez, eu teria que listar
todos os países e climas que deram ao mundo uma camaradagem tão bonita
entre homens e mulheres, mas o espaço não permite. Só estou levantando
este ponto para enfatizar o absurdo dos argumentos do autor de ``Louise
Michel''.

Louise Michel odiava os babados das roupas femininas e todos os demais
requisitos que compõem a fragilidade da mulher numa sociedade
pervertida, além disso era também descuidada e desorganizada. No que diz
respeito ao suposto primeiro argumento, devo fazer um esclarecimento a
von Levetzow. É verdade que muitas mulheres que se emanciparam da
vergonha do passado desenvolveram uma nova concepção de beleza e
elegância ao se vestir. Mas, na maioria dos casos, tudo isso foi apenas
um protesto contra os trapos, o desperdício e a estupidez de toda a
parafernália que compunha as vestimentas das mulheres comuns. De um
ponto de vista científico, sociológico e moral, essas mulheres
insistiram que a marca da escravidão do seu sexo eram suas roupas e que
elas não poderiam ser realmente livres a menos que transvalorizassem o
valor das coisas que promoviam a sua escravidão. Todas essas mulheres
são homossexuais? Não mais do que Louise o foi. Louise, que dedicou sua
vida à causa da humanidade, que não apenas se engajou na luta pela
existência de sua mãe e de si mesma, mas, acima de tudo, no movimento
que absorveu a maior parte dos seus pensamentos e toda a sua
energia.\footnote{Filha ilegítima e de família humilde, Michel era muita
  próxima da mãe que tinha uma saúde frágil. Mesmo na prisão Michel
  continuou a lhe dar suporte, obtendo, em 1884, autorização para ficar
  ao lado da mãe nos seus últimos dias de vida.} Deveria ela, então, ter
passado horas na frente do espelho, explorando costureiras e torturando
vendedoras em nome da busca vã pelos modelos mais recentes? Deve ela,
portanto, ser considerada uma uraniana, porque se vestiu com sensatez e
prestou pouca atenção ao que é comumente considerado a beleza da mulher?
Na verdade, se o autor não possui provas melhores do que essa sua
alegação, ele deveria ter se abstido de decifrar o seu caso.

Louise Michel possuía um corpo retangular e tinha traços masculinos. Não
é verdade, como diz o autor, que ela tinha uma voz masculina. Eu a ouvi
falar, quando ela tinha 66 anos de idade; sua voz era de contralto,
bela, profunda e melódica e atingia diretamente o coração dos seus
ouvintes. Havia uma notável simplicidade que explica o grande poder que
ela tinha sobre o seu público. Quanto ao seu rosto, é bem óbvio para mim
que von Levetzow nunca teve a oportunidade de ver Louise sorrir. Caso
tivesse tido, não teria visto nela, um homem. O efeito iluminador que o
sorriso e os belos olhos de Louise Michel criavam, era uma unanimidade
entre todos aqueles que eram próximos a ela. Este argumento, também,
parece muito fraco.

Aqui, chegamos à afirmação mais importante de von Levetzow. Segundo o
relato do autor, Louise recusou dois pretendentes por não se sentir
atraída por homens. É bastante significativo, porém, que isso tenha
acontecido aos 12 e 13 anos de idade e que ambos os homens, seus
pretendentes, tinham idade suficiente para serem seu pai; acrescente-se
a isso o fato de que na verdade eles foram comprá-la. Ela mesma expressa
sua indignação e repulsa contra esses homens, na página 330. Também a
sua atitude em relação à instituição do casamento é muito significativa
(página 332).\footnote{Goldman está se referindo aqui a trechos da
  autobiografia de Louise Michel, \emph{Mémoires de Louise Michel}
  (1886), citados no ensaio de von Levetzow. A paginação a que ela
  remete o leitor diz respeito, portanto, ao artigo de von Levetzow
  (\emph{Jahrbuch für sexuelle Zwischenstufen}, vol. VII, 1905 pp.
  307-370).} Louise se ressentia do casamento sem amor, assim como
deveria toda mulher que se preze. Não obstante, em nenhuma passagem dos
seus escritos, ela expressou oposição ao amor sem casamento. Tampouco
ela escreveu sobre a necessidade de tornar pública uma experiência tão
profundamente íntima e privada quanto a vida amorosa entre duas pessoas.
Abaixo segue a prova disso:

Louise Michel teve uma experiência amorosa com um professor quando era
muito jovem e trabalhava, também ela, como professora. Mais tarde, após
retornar da Nova Caledônia, ela viveu por um tempo com um camarada
belga. E se não teve mais experiências desse tipo, foi, provavelmente,
porque, conforme ela mesma declarou: ``eu dei o meu coração à
Revolução''. Sim, esse era o amor de Louise. Ao longo de toda a sua
vida, ela foi dedicada a esse amor. Tipos como Louise Michel não podem
ter um amor pessoal, que de algum modo venha a interferir na grande
paixão por um ideal, e a paixão de Louise pelo seu ideal foi o elemento
mais poderoso em sua vida, acompanhando-a até o túmulo.

É verdade que ela tinha várias amigas que amava, mas não da maneira que
von Levetzow supôs. Muitas mulheres modernas que não têm muito romance a
permear sua vida pessoal, dedicam-se à camaradagem e à amizade com
pessoas do seu sexo, assim como acontece com muitos dos grandes homens.
A razão para isso, no caso das mulheres, é que elas encontram uma
compreensão mais profunda entre os membros de seu próprio sexo do que
com os homens do seu tempo. O problema é que o homem moderno ainda se
assemelha demais ao seu antepassado Adão, não diferindo muito, na sua
atitude em relação à mulher, de um homem mediano qualquer. Por outro
lado, a mulher moderna já não se satisfaz com um homem que seja tão
somente seu amante; ela quer compreensão, camaradagem, quer ser tratada
como um ser humano, e não como um objeto de gratificação sexual. Uma vez
que ela nem sempre pode encontrar isso no homem, ela se volta para suas
irmãs. É precisamente \emph{porque não há o elemento sexual} entre elas
que elas podem compreender melhor uma a outra. Em outras palavras, ao
invés de se sentir atraída por suas amigas mulheres por conta de
tendências homossexuais, Louise se atraía por elas justamente porque era
uma mulher e precisava da companhia de mulheres. Há mais uma coisa que
era muito proeminente em Louise Michel: o seu instinto materno.
Apaixonada por crianças, ela era extremamente maternal com todas as
crianças abandonas que conseguia ajudar; foi o amor materno que a levou
a adotar Charlotte Vauzelle, para educá-la e compartilhar com ela os
seus escassos proventos.\footnote{A adoção que Goldman se refere se deu
  no sentido emocional e intelectual (e, segundo o que coloca aqui,
  também financeiro), mas não legal.} Charlotte nunca foi, nem no
sentido espiritual, a namorada de Louise Michel. O fato é que Louise
pagou caro por sua dedicação a Charlotte. Juntamente com o irmão,
Charlotte tornou os últimos anos de vida de Louise Michel miseráveis. De
certo modo, eles a mantiveram prisioneira, não a deixavam em paz fosse
para receber seus amigos ou para ir morar com eles. Charlotte violava a
correspondência de Louise e a observava atentamente. A razão de tudo
isso é que Charlotte e seu irmão viviam às custas de Louise Michel e
tinham um medo mortal de que mais alguém pudesse se beneficiar da sua
grande generosidade. Qualquer que fosse o caso, é ridículo considerar
que Charlotte Vauzelle era a namorada de Louise Michel.

O amor de Louise Michel pela escultura e sua apreciação pela música de
Wagner são apresentados como evidência de traços homossexuais. Von
Levetzow afirma que nenhuma mulher é capaz de criatividade e gosto
artístico e musical. Generosamente, ele chega a admitir que Francisco
Holmés, uma mulher francesa de origem escandinava, era uma grande
compositora.\footnote{Provavelmente, Goldman está se referindo aqui a
  Augusta Mary Anne Holmès (1847-1903), compositora e pianista francesa
  de ascendência irlandesa.} Seja como for, ele se apressa em
acrescentar que, de acordo com a sua fotografia, ela parece masculina.
Não acho que valha a pena entrar nesse argumento. O fato de haja poucas
mulheres que sejam grandes compositoras não torna as que são,
homossexuais. Eles são simplesmente pioneiras em um domínio que até
então não foi explorado por muitas mulheres. Quanto ao amor por Wagner,
a verdade é que, em geral, as mulheres ouvem a música wagneriana e a
entendem muito mais e melhor do que homens. Talvez isso se dê dessa
maneira, porque o espírito ilimitado, elementar à música de Wagner,
afete as mulheres como uma força liberadora das emoções reprimidas e
ocultas nas suas almas. Não parece ser necessário enfatizar que as
mulheres não são apenas capazes de apreciar a escultura, mas que há um
grande número de mulheres escultoras de mérito não desprezível.

Há, porém, uma coisa em que concordo plenamente com von Levetzow: é
quando ele afirma que Louise Michel estava tão intrinsecamente envolvida
com o anarquismo que, para compreender sua personalidade e sua natureza
complexa, é preciso também adentrar na discussão da sua filosofia
social. Mas, como ele mesmo diz, o \emph{Jahrbuch} não é lugar para
isso. Ainda que fosse, não creio que o autor estaria em posição de fazer
uma análise do anarquismo, já que ele parece não saber absolutamente
nada sobre o assunto. Afinal, se não for isso, como é possível entender
a sua interpretação contida na página 315 (linhas 4 a 7)? Que insulto à
memória de Louise Michel e à inteligência dos leitores do
\emph{Jahrbuch}!\footnote{Ao especular sobre quem seria o pai de Louise
  Michel, von Levetzow chega à conclusão que ele seria integrante de
  alguma família proprietária de terras de antiga linhagem feudal. Daí
  formula a passagem aqui execrada por Goldman: segundo ele, essa
  ascendência não é, de modo algum irrelevante, ``porque a história do
  anarquismo parece sugerir que pelo menos uma boa porcentagem de sangue
  tirano é necessária para produzir uma vila anarquista''.}

Conheço Louise Michel já há muitos anos. E muito antes de conhecê-la,
conhecia as suas ideias e o preço que ela pagou ao lutar por elas. A sua
força, o seu martírio e, mais do que isso, o seu amor ilimitado pela
humanidade eram para mim como uma chama que me purificava e iluminava.
Eu a conheci pela primeira vez em 1896, em Londres; na época,
encontramo-nos com frequência e foi ela quem me ensinou sobre a história
da luta heroica da Comuna de Paris. Louise nunca falou de si mesma e
sobre o papel que ela desempenhou nessa luta.

Eu estive com ela novamente em 1899, em Londres, e depois, em 1900,
quando pela primeira vez, depois de muitos anos, Louise Michel retornou
a Paris. Foi durante esses dois períodos que eu tive a oportunidade de
conviver bastante com ela e ouvir ela contar alguns episódios da sua
vida, já que era minha intenção escrever uma biografia sua. Mas ela era
tão excessivamente reticente em relação a tudo que dissesse respeito a
si mesma que relutava em discutir a sua vida. Por outro lado, quando
falava dos outros, ela sempre ficava radiante, era como se o seu rosto
se iluminasse com um brilho divino: fosse sobre os seus camaradas, quem
ela auxiliou e cuidou enquanto estava na Nova Caledônia; ou, mesmo sobre
as criaturas irracionais. Dentre as características de Louise Michel
estava a sua grande simpatia pelos animais. A casinha em que ela vivia
em Londres era como que uma coleção perfeita de gatos e cães abandonados
que ela recolhia à noite, quando a caminho de casa. O seu amor por
gatos, em especial, era certamente qualquer coisa, menos masculino. É
verdade que von Levetzow relata o fato de que Louise, quando nas
barricadas e rodeada por balas, resgatou um gato que estava paralisado
pelo medo. A história não mencionou ainda nenhum homem que, em perigo,
tenha feito uma coisa dessas. Não quero dizer que um homem não se poria
em risco para salvar uma criança ou até um cachorro; certamente, porém,
nunca um gato.

A supostamente masculina Louise Michel, que era desorganizada e incapaz
de se arrumar, enfim, que não era uma doméstica, ainda assim aprendeu a
tricotar, costurar, lavar e cozinhar para seus companheiros exilados na
Nova Caledônia, além de cuidar deles quando doentes com toda a ternura
do seu grande coração de mãe, e de dar suporte aos seus espíritos quando
o que havia de mais terrível em suas vidas aconteceu.\footnote{Após a
  repressão sangrenta que deu fim à Comuna de Paris, Louise Michel e
  muitos dos seus companheiros, aproximadamente 4500, foram exilados na
  Nova Caledônia.}

Lembro-me de uma noite maravilhosa em Paris. Amigos anarquistas
organizaram um pequeno jantar para Louise, para o qual fui convidada.
Louise, vestida de preto como habitual, com apenas uma gola de renda e
punhos da manga brancos para dar algum alívio, tinha o rosto tão corado
quanto as rosas sobre a mesa, o qual, por sua vez, era como que
emoldurado pelos seus cabelos encaracolados, cor de prata. Ela estava
radiante de alegria por estar de volta à cidade dos seus sonhos e lutas,
cercada por seus amigos íntimos. Ela estava mais falante do que em
qualquer outro momento em que antes eu a ouvira, mais disposta a nos
deixar olhar o interior da sua alma. Em nenhum momento, Louise
demonstrou, sequer remotamente, alguma característica masculina, como
tampouco tendências homossexuais. Estou certa de que tais indícios não
teriam me escapado se, de fato, estivessem presentes em Louise Michel,
posto que, como eu disse no começo do meu artigo, fiz um estudo sobre
homossexualidade a partir do que há de melhor literatura, conheço muitos
homossexuais e facilmente detecto inclinações homossexuais nas pessoas.
Não havia vestígios de homossexualidade em Louise Michel.

Entre os amigos de Louise Michel estavam os maiores homens do seu tempo:
Piort Kropotkin, Malatesta, Elisée Reclus, Malato, Rocker. Alguns deles
moravam bem perto dela, em contato quase diário. Houvesse o menor
indício de homossexualidade, eles teriam percebido; certamente isso
seria do conhecimento dos seus camaradas. Recentemente, falei com meu
amigo e camarada, Rudolf Rocker, sobre esse aspecto tratado no artigo de
von Levetzow. Ele também me garantiu que nunca, em nenhum momento,
nenhum dos amigos íntimos de Louise Michel viu o menor indício de
inclinações homossexuais. Diga-se de passagem, inclusive, que Rudolf
Rocker, assim como eu, está absolutamente livre de qualquer preconceito
em relação aos homossexuais. Nosso único desejo é retratar Louise Michel
tal como ela realmente era -- uma mulher excepcional, dotada de grande
intelecto e espírito maravilhoso. Ela representava o novo tipo de
feminilidade, ainda que tão antigo quanto a raça, tão sábio quanto o
tempo e dotado de uma alma repleta de amor pela humanidade. Em suma, uma
mulher completa, livre do preconceito e dos costumes que, por séculos,
mantiveram a mulher em cativeiro, condenando-a à posição de escrava
doméstica e sexual. Com Louise Michel surge a nova mulher que é capaz
dos atos mais heroicos, ao mesmo tempo em que continua sendo uma mulher
nas suas paixões e vida amorosa.

Caro Dr. Hirschfeld, tentei, da maneira mais concisa possível, analisar
criticamente as alegações de Frhr. von Levetzow. Tenho certeza que
concordará comigo que nem a questão da homossexualidade, como tampouco
os homossexuais ganharão qualquer coisa por meio da distorção dos fatos.
Foi por esse motivo que me comprometi a refutar os erros do artigo e
jamais por qualquer outro. Espero que considere minhas críticas
convincentes e que não apenas publique a minha réplica, como já
gentilmente me ofereceu, mas que também retire a fotografia de Louise
Michel da sua galeria de uranianos.

Atenciosamente,

Emma Goldman

\chapter[Mulheres heróicas da Revolução Russa (1925)]{Mulheres heróicas da Revolução Russa\footnote[*]{Texto não publicado em
  vida, a datação de 1925 é aproximada.}}



A Rússia pré-revolucionária tornou-se única na história mundial pelo
grande número de mulheres excepcionais e heroicas que contribuíram para
o movimento da sua libertação. Teve início com as esposas dos
``decembristas'', os primeiros rebeldes políticos contra o czarismo
autocrático, quase um século atrás, que acompanharam voluntariamente os
seus maridos no exílio na Sibéria, até o último dia da dinastia dos
Romanov.\footnote{Ao optarem por seguir seus maridos no exílio à
  Sibéria, as esposas dos decembristas tiveram de abdicar dos títulos de
  nobreza, propriedades, além de não poderem levar consigo os seus
  filhos. Os que viessem a nascer no exílio não poderiam ter o nome do
  pai, sendo, portanto, ilegítimos e o seu retorno só seria permitido
  após a morte do marido. Na Sibéria, desenvolveram um intenso trabalho
  de apoio aos presos políticos e na educação dos jovens e crianças das
  aldeias.} As mulheres russas participaram de todas as formas de
atividade revolucionária e foram em direção à morte e à prisão com um
sorriso nos lábios.

No poema vívido e pujante, ``Mulheres russas'', o poeta Nekrassóv pagou
um alto tributo à força e ao valor das mulheres que sacrificaram a sua
riqueza, posição social e cultura para seguir o caminho árduo ao longo
das planícies congeladas do norte, a fim de compartilhar o destino cruel
dos seus maridos presos e exilados. Após Nekrassóv, foi a vez de Ivan
Turguêniev que, com sentimentos elevados e admiração empática, pintou a
imagem das mulheres russas revolucionárias do seu tempo. No seu soberbo
poema em prosa, intitulado ``No limiar'', ele imortalizou as idealistas
ardentes, as mulheres russas do tipo de Sofia Perovskaia, cuja fé
apaixonada e devoção altruísta à liberdade iluminaram, tal qual um
farol, o horizonte sombrio da Rússia do início dos anos
oitenta.\footnote{Sofia Perovskaia (1853-1881) era um dos membros mais
  destacados do relatiy8conhecido grupo terrorista russo Naródnaia Vólia
  {[}``Vontade do povo''{]}, responsável pelo assassinato do czar
  Alexandre II, em 1881. Foi enforcada em praça pública, junto a outros
  companheiros no mesmo ano. Foi a primeira mulher a ser executada na
  Rússia moderna.}

A Revolução de Fevereiro de 1917 abriu as portas da prisão para os
sobreviventes da tortura, do calabouço e do exílio siberiano
distribuídos pelo sistema czarista aos seus oponentes políticos. Em
triunfo, foram trazidos à Moscou e Petrogrado, dezenas de
revolucionários das gerações mais jovens, dentre os quais encontravam-se
nomes reverenciados como os de Maria Spiridonova, sua amiga íntima
Alexandra Izmailovitch, Irena Kakhovskaia, Evgenia Ratner, Olga Taratuta
-- representantes de vertentes políticas várias, mas todas inspiradas no
amor ao povo e na devoção a sua causa.

Olga Taratuta, filha de intelectuais, embora dotada de físico delicado,
possuía uma mente poderosa e foi, em certo sentido, uma pioneira. Quando
tinha apenas vinte anos, organizou, com vários amigos, o primeiro grupo
anarquista do sul da Rússia. Era um empreendimento perigoso, e suas
atividades logo atraíram a atenção da polícia política. Presa no início
da revolução de 1905, Olga foi condenada a 30 anos de katorga (trabalho
forçado na prisão), em Odessa. Engenhosa e ousada, conseguiu escapar, e
retomou o seu trabalho anterior, dessa vez, com um nome falso. Durante
um tempo considerável todos os esforços da polícia para encontrá-la não
renderam frutos, mas em 1906 seu disfarce foi descoberto: ela foi presa
novamente e sentenciada, mais uma vez, a 30 anos de prisão. Quando
retornou à liberdade, em 1917, Olga se dedicou ao trabalho político na
Cruz Vermelha, ajudando as vítimas do regime do comandante militar
Skoropadsky,\footnote{Com a revolução de fevereiro de 1917, o general
  conservador Pavlo Skoropadski organizou uma reação na Ucrânia,
  atribuindo-se o título de Hetman da Ucrânia, o qual ele pretendia
  hereditário. Não se manteve no poder, porém, para além do ano de 1918.}
na Ucrânia, e, posteriormente, dedicou-se aos cuidados dos novos grupos
de presos políticos criados pelo Estado Comunista.

Nos últimos anos dos 1920, uma conferência com anarquistas de toda
Rússia, estava sendo preparada em Cracóvia. Embora esse encontro
estivesse sendo organizado com o conhecimento e o consentimento do
governo soviético, todos os delegados foram presos na véspera da
conferência, sem aviso ou explicação. Em meio às centenas de
prisioneiros estava também Olga Taratuta. Ela foi mandada para a prisão
de Boutyrki, em Moscou, o mesmo presídio em que tantos dos seus
companheiros haviam sofrido e morrido nos dias do regime Romanov. Lá,
Olga foi obrigada a vivenciar a experiência mais dolorosa de sua vida
cheia de acontecimentos notáveis. Na noite do dia 25 de abril, a ala dos
presos políticos foi invadida pela Tcheka,\footnote{Tcheka é o acrônimo
  da ``Comissão Especial Pan-Russa de Luta contra a Contrarrevolução e a
  Sabotagem'', primeira forma de organização da polícia política secreta
  soviética.} os prisioneiros foram atacados enquanto dormiam e, após
uma série de maus tratos, foram arrastados para a estação ferroviária --
com alguns deles, vestindo nada além das suas roupas noturnas -- e
transferidos para outros presídios.

Olga se viu, então, encarcerada na temida prisão de Orlov, que, sob o
governo de Nicolau II, era o local central para a ``distribuição'' dos
presos. A administração e o regime daquela prisão eram tais que,
rapidamente, levaram os presos políticos a uma greve de fome como
protesto contra o tratamento que estavam recebendo. Olga foi novamente
removida para outra prisão, e de lá enviada para o exílio na região
sombria de Veliki Ústiug, sendo, por fim, enviada para a prisão de Kiev,
a mesma onde anteriormente havia se dedicado, tão diligentemente, aos
cuidados dos comunistas presos na reação liderada por Skoropadski. Uma
carta recente de Olga a um amigo que vive no exterior contém uma
observação significativa: a de que a perseguição imposta pelo governo
soviético lhe roubou mais vitalidade do que todos os anos de
encarceramento que ela sofreu nas mãos da autocracia Romanov.

Diferentemente de Olga Taratuta, a maioria das outras heroínas da
revolução russa tem origem proletária. Dentre elas, Leah Gotman e Fanya
Baron são duas mulheres anarquistas de personalidade excepcional. Na sua
adolescência, ambas partiram da Rússia em direção à América, onde foram
empregadas nas fábricas e participaram ativamente do movimento dos
trabalhadores. Eu conheci bem essas garotas, espécimes esplêndidos da
independência feminina, belas, dotadas de sentimentos elevados e
intelecto forte. À primeira chamada da Revolução de Fevereiro, essas
duas garotas, como muitos refugiados russos, correram para a sua terra
natal. E, assim, ajudaram a realizar a Revolução de Outubro. Leah e
Fanya sentiram que o seu lugar era junto ao proletariado, trabalharam,
particularmente, com os mujiques do sul, em meio aos agricultores da
Ucrânia, a quem deram todo o amor e devoção inerentes às suas naturezas
tão ricas. Subsequentemente, as duas garotas desenvolveram uma série de
atividades culturais com os camponeses rebeldes liderados pelo famoso
Bat'ka (``Pequeno Pai''), Nestor Makhno.\footnote{Nestor Makhno
  (1888-1934) foi um dos principais nomes do anarquismo na Ucrânia. Como
  muitos revolucionários presos pelo regime czarista, foi libertado na
  ocasião da Revolução de Fevereiro. Lutou ao lado dos bolcheviques, com
  o exército independente do qual era comandante, o chamado Exército
  Negro, na guerra civil russa que se seguiu à revolução. As relações
  tensas entre os bolcheviques e os anarquistas makhonistas, como
  ficaram conhecidos os membros do movimento liderado por ele,
  concluiu-se, em 1920, com a invasão do território makhonista pelo
  Exército Vermelho, que estendeu a sua perseguição implacável aos
  anarquistas como um todo.}

As mãos do Kremelin uma vez erguidas contra Makhno, também caíram
pesadamente sobre Leah Gotman e Fanya Baron. Ambas foram presas na
véspera da conferência em Cracóvia, mencionada acima, e mandadas para a
prisão de Boutyrki, onde também foram vítimas do ataque da Tcheka, na
noite de 25 de abril de 1920. Arrancada da cama na calada da noite, Leah
foi arrastada nas escadas pelos cabelos e forçada a permanecer por
horas, semivestida, no pátio da prisão junto a outros presos políticos,
enquanto esperava ser transferida a algum destino desconhecido. Ela
permanece presa desde então, sendo agora uma das desafortunadas detentas
do terrível Mosteiro de Solovetsky, situado na zona do Ártico.

Fanya Baron, que sempre me impressionou pela sua coragem ilimitada e
espírito excepcionalmente generoso, pertence ao tipo raro de mulher que
pode realizar as mais difíceis tarefas da prática revolucionária com
calma, graça e altruísmo absoluto. Após o ataque em Boutyrki, ela foi
transferida para a prisão de Riazã, de onde logo escapou, voltando a pé,
para Moscou, sem nenhuma ajuda. Chegou lá sem dinheiro e praticamente
sem roupas, sendo obrigada, por conta dessa condição desesperadora, a
buscar refúgio com o irmão do marido, onde foi descoberta pela Tcheka.
Essa mulher de grande coração, que serviu à causa da Revolução por toda
a vida, foi morta pelo Partido que finge ser a guarda avançada da
Revolução. Não satisfeitos com o assassinato de Fanya Baron (em setembro
de 1921), os comunistas mancharam com o estigma de ``banditismo'' a
memória de sua vítima assassinada.

Não apenas os anarquistas, mas todos os membros de qualquer outro grupo
político, com os social-revolucionários de direita e de esquerda, os
mencheviques, os maximalistas e mesmo os comunistas, tiveram de pagar
caro à brutalidade autocracia comunista. A partir de agora, citarei
algumas das personalidades mais extraordinárias.

Evgenia Ratner, uma jovem de mente afiada e personalidade forte,
ingressou no Partido Socialista-Revolucionário, logo após completar seus
estudos em medicina na Suíça. Suas atividades, depois que retornou à
Rússia, envolveram-na repetidamente em dificuldades para com as
autoridades czaristas, até que foi condenada a muitos anos de prisão.
Libertada pela Revolução de Fevereiro de 1917, sua habilidade e força
excepcionais a levaram a ser eleita como membro do Comitê Central do seu
partido, e, ao mesmo tempo, também foi escolhida pelo campesinato como
um de seus representantes no Soviete de Moscou. Logo após seu partido
ser banido pelos bolcheviques, Evgenia foi presa em 1919 e julgada, em
1922, juntamente com onze camaradas. Todos foram condenados à morte.

A intercedência do mundo ocidental, com um protesto internacional de
grande fôlego contra a execução da sentença -- apoiado por homens como
Anatole France, Romain Rolland e outros --, salvou as vidas dos doze
social-revolucionários, Evgenia Ratner entre eles. Ela agora leva uma
vida miserável na prisão de Boutyrki.

Dos socialistas-revolucionários de esquerda, Irina Kakhovskaia,
Alexandra Izmailovitch e Maria Spiridonova sofreram o maior martírio.
Kakhovskaia, neta do general Kakhovski, o famoso rebelde ``decembrista''
que se opôs à autocracia de Nicolau I, é uma mulher reconhecida pelos
seus dotes literários e idealismo revolucionário. Ela começou seu
trabalho nos movimentos de libertação da Rússia, quando era muito jovem,
em 1904. Não demorou muito, foi presa e sentenciada a 20 anos de
katorga, sendo, posteriormente, transferida para Akatuy, um dos lugares
mais temidos do exílio czarista. Em 1914, foi autorizada a viver, em
exílio, na região da Transbaikalia, do qual foi libertada pela Revolução
de Fevereiro de 1917.

Após o seu retorno do exílio, Irina Kakhovskaia tornou-se uma das
trabalhadoras mais valiosas do Partido Socialista-Revolucionário de
Esquerda, era muito respeitada pela sua compreensão de psicologia
camponesa e das necessidades do proletariado. Após o Tratado de
Brest-Litovsk e da ocupação alemã da Ucrânia, as autoridades alemãs
prenderam Irina como participante da conspiração contra a vida do
general Eichorn, governador militar da Ucrânia, morto por um
revolucionário do seu partido, B. Donskoy. Kakhovskaia foi torturada e
condenado à morte. Felizmente, a eclosão da revolução na Alemanha
impediu a sua execução e ela foi salva.

Irina continuou a trabalhar nas suas convicções políticas, mas, em 1921,
ela foi presa novamente, desta vez pelos bolcheviques, por quem foi
exilada em Kaluga, na Sibéria.

Enquanto estava na prisão, Irina Kakhovskaia escreveu as suas
interessantíssimas memórias, a história incomum de uma personalidade
muito singular. Romain Rolland, depois de examinar o seu trabalho,
disse: ``Sou contrário às ideias de Kakhovskaia, mas sua narrativa tem
uma qualidade humana, ou antes, sobrehumana, muito cativante. É um
documento psicológico do valor mais elevado que pode haver. A
simplicidade absoluta da narradora, sua habilidade, genuinamente russa,
de ver de modo objetivo, sua energia inacreditável inteiramente dedicada
à causa que tem no coração -- tudo isso desperta profunda admiração no
leitor, independentemente de qual seja a sua posição em relação ao valor
da ação realizada ou almejada. Que heroísmo, paciência, completa
abnegação, que tesouros da alma, a humanidade desperdiça com propósitos
terríveis e vergonhosos''.

Alexandra Izmailovitch, filha de um general do exército russo, é outro
exemplo de uma jovem que a autocracia dos Romanov, impeliu a atos
individuais de violência como única forma de protesto possível sob um
regime despótico. Em 1906, ela atentou contra a vida do governador
Kurlov da província de Minsk, responsável pela maior parte dos pogroms
diabólicos contra os judeus. Condenada à Sibéria por toda a vida, ela
foi libertada com os outros ativistas políticos, em 1917. Na condição de
membro do Partido Socialista-Revolucionário de Esquerda, tornou-se uma
figura de liderança no Soviete de Deputados Camponeses de Toda Rússia.
Quando os bolcheviques decidiram ``liquidar'' seu Partido ``para
sempre'', em 1919, ela foi presa juntamente com vários de seus
companheiros, permanecendo, quase que ininterruptamente, na prisão desde
então.

O aspecto mais característico dessa mulher extraordinariamente dotada e
enérgica é a devoção, ao longo de toda a sua vida, à sua amiga e
camarada Maria Spiridonova. Elas passaram onze anos juntas na Sibéria,
retornaram à Rússia para unir seus esforços na libertação do povo, e
foram presas pelo governo bolchevique. Desde então, já há muitos anos,
estão presas. Não é exagero dizer que o cuidado e a devoção de Alexandra
Izmailovitch à amiga são as principais causas de Maria Spiridonova ainda
estar entre os vivos.

Maria Spiridonova é, sem dúvida, uma das figuras mais notáveis ​​e
heroicas do movimento revolucionário russo dos últimos vinte anos. De
família aristocrática, bela e culta, a jovem Maria abandonou o luxo e
sua posição social para se dedicar à causa dos oprimidos.
Bem-intencionada e empática, ela não pôde suportar sem protestar a
injustiça e tirania que testemunhava de todos os lados. Com 18 anos,
atentou contra a vida do general Lukhanovsky, governador da província de
Tambov, que foi universalmente execrado por sua selvageria asiática
contra o campesinato.

Os czares russos nunca foram parciais no tratamento às ativistas
políticas: eram igualmente implacáveis ​​com todos os seus oponentes,
fossem homens ou mulheres. Mas no caso de Maria Spiridonova, os capangas
de Nicolau II superaram até os métodos de Ivan, o Terrível. Após ser
presa, Maria foi espancada até o ponto da insensibilidade, suas roupas
literalmente arrancadas do corpo, e a jovem foi entregue a uma guarda
embriagada que encontrou diversão em queimar sua carne nua com cigarros
acesos. Depois de semanas à beira da morte, Maria foi condenada à
execução.

A tortura de Spiridonova comoveu todo o mundo ocidental, cujos protestos
a salvaram do cadafalso. Ela foi ``perdoada'', tendo a sua pena comutada
por prisão perpétua na Sibéria. Os efeitos da sua experiência terrível
provocaram lesões nos pulmões, aleijamento numa das mãos e perda da
visão em um olho. Não obstante, ainda que fisicamente debilitada e
destroçada, seu espírito permaneceu em chamas.

Poucos ativistas políticos receberam uma ovação popular tão unanime, da
Sibéria a Petrogrado e Moscou, como Maria Spiridonova, após sua
libertação da prisão em 1917. Ela, porém, não desperdiçou um único
segundo com o gozo da sua recém-conquistada liberdade. Dedicou-se ao
trabalho com todo o ardor característico à sua personalidade intensa,
organizando os camponeses, inspirando e dirigindo as energias então
despertas do povo russo. Tornou-se a líder adorada dos vários milhões de
camponeses da Rússia, a alma de todas as suas aspirações de longa data e
a porta-voz das suas necessidades e esperanças. Como a figura de maior
destaque do Partido Socialista-Revolucionário de Esquerda, Maria exerceu
uma influência tremenda no soviete dos camponeses, elaborou um plano
abrangente para a socialização da terra, o problema mais vital da vida
russa.

Já em 1918, Maria Spiridonova percebeu que a Revolução corria grande
risco, e da parte de alguns de seus supostos amigos, ao invés dos
inimigos. Ela viu a crescente autocracia do Estado comunista e se
colocou contra ela. A ruptura final entre o seu partido e os
bolcheviques veio do Tratado de Brest-Litovsk, que Spiridonova condenou
tanto por questões de princípio, quanto por razões práticas. Pouco após
isso, ela foi presa junto com 500 delegados no Congresso dos camponeses.

Quando cheguei à Rússia, os bolcheviques me disseram que Maria
Spiridonova havia sofrido um colapso nervoso e que, por conta disso,
teria sido alocada num sanatório, para que recebesse os melhores
cuidados. Mas logo descobri que Maria havia escapado desses ``melhores
cuidados'' e estava vivendo em Moscou disfarçada de camponesa, como
costumava fazer nos dias de czarismo. Recentemente, a fortuna me
favoreceu e tive a oportunidade de passar vários dias com essa mulher
extraordinária. Não encontrei nela nenhum vestígio de histeria -- na
verdade, o seu equilíbrio, ponderação e objetividade ao narrar os
eventos que marcaram o seu retorno à Rússia são absolutamente
admiráveis.

Alguns meses depois, no outono de 1920, a Tcheka retornou, mais uma vez,
à sua atividade de caçar conspirações. Numa das suas numerosas incursões
em Moscou, encontraram Maria Spiridonova, doente de tifo. Ela foi presa
e removida para Ossoby Otdel -- a Seção Secreta da Tcheka. Em 1921,
quando Maria estava quase à beira da morte, os esforços dos seus amigos
garantiram a concessão da sua libertação temporária -- sob a condição de
retornar à prisão assim que a sua saúde estivesse melhor. A única
alternativa era deixar Maria morrer na prisão por negligência ou
devolvê-la -- quando melhor de saúde -- para os ``melhores cuidados''.
De fato, assim que ela começou a se recuperar, a Tcheka se apoderou dela
novamente. Guardas com cães farejadores foram colocados na casa onde
Spiridonova estava sendo assistida pela sua devotada amiga, Alexandra
Izmailovitch. Todos os seus passos eram observados e a existência
tornou-se tão insuportável que a torturada terminou por exigir ser
levada de volta à prisão. Juntamente com a sua inseparável Izmailovitch,
ela foi enviada para um dos locais mais distantes e inóspitos da
província de Moscou, de onde vêm agora as mais tristes notícias, como a
de que Spiridonova foi obrigada a recorrer ao método desesperado da
greve a fome como protesto contra a implacável perseguição. De fontes
confiáveis, ​​acaba de chegar a informação de que tanto Izmailovitch
quanto Spiridonova estão exiladas nas regiões mais selvagens do
Turquestão.\footnote{Em 11 de setembro de 1941, Spiridonova e
  Izmailovitch, em conjunto com mais de 150 presos políticos, foram
  executadas via ordem pessoal de Stalin, na ocasião da invasão das
  tropas nazistas na URSS. Essa execução em massa ficou conhecida como
  ``Massacre da Floresta Medvedev''.}

O martírio das mulheres heroicas da Rússia tornou-se mais doloroso e
intenso sob a tirania da ditadura bolchevique do que quando nos dias do
czarismo. Naqueles tempos, o sofrimento era meramente físico, nada
poderia afetar o espírito. Elas sabiam que enquanto eram odiadas pela
autocracia, desfrutavam do respeito e amor das vastas massas do povo
russo. Na verdade, o ``povo simples'' as tomava como ``santas'' que
sofriam por sua causa; no czarismo, a influência moral exercida pelos
ativistas políticos nas prisões, katorgas e exílios era muito grande.
Mas tudo isso mudou agora. Os novos autocratas da Rússia desacreditaram
os ideais do socialismo e macularam a honra dos seus grandes expoentes.
Não há opinião pública na Rússia que não seja a do partido governante, e
os mártires -- homens e mulheres -- da Rússia revolucionária se
transformaram em párias, no sentido mais amplo que pode haver. Eles não
têm nada para os compensar, não podem sequer apelar à consciência do seu
país, pois ela foi paralisada. Inclusive, não apenas a consciência da
Rússia, mas a consciência do mundo como um todo parece silenciada.

O que aconteceu com o sentido de justiça e solidariedade que
anteriormente se estendia do mundo ocidental às vítimas políticas do
regime czarista? Naquele tempo, homens e mulheres da Inglaterra, de modo
corajoso, denunciavam em seus protestos as iniquidades cometidas na
Rússia e eram solidários com os que eram perseguidos politicamente por
conta das suas opiniões. Agora, diante das evidências esmagadoras da
opressão e perseguição extremamente cruéis que ocorrem na Rússia, o
mundo está silencioso e insensível. Os mártires heroicos estão
abandonados à misericórdia da Tcheka, para sofrer o Gólgota do corpo e
do espírito, em nome de um ideal que há muito tempo foi traído pelo
Estado Comunista com a sua ditadura do Partido.

\chapter[As visões de Emma sobre o amor (1926)]{As visões de Emma sobre o amor\footnote[*]{Texto originalmente publicado no
  \emph{NEA Service}. O título e as subseções, que se optou por manter
  aqui, foram dados pelo jornal; no rascunho o texto curto, sem
  subseções é intitulado: ``Minha atitude em relação ao casamento''} (1926).}

Nota do jornal -- Nos tempos modernos, nenhuma mulher posicionou-se, de
modo mais abertamente hostil, contra o casamento do que Emma Goldman,
que, há oito anos, estava entre os deportados enviados de volta à
Rússia. Ela fez uma série de palestras na América, em que tratou sobre o
tema do casamento e escreveu panfletos. De repente, veio de Montreal
(Canadá) a notícia de que Emma Goldman havia chegado lá sob o nome de
Sra. E. Colton. Após uma vida de antagonismos, Emma Goldman se casou com
um trabalhador das minas. Seus seguidores ficaram surpresos e ainda
estão. O que isso poderia significar? A NEA enviou um representante até
ela para saber de primeira mão o que Emma Goldman agora pensa sobre o
casamento. Aqui, apresentamos um artigo exclusivo, escrito pela próprio
Goldman, o segundo de uma série de cinco, superando suas opiniões
hoje.\footnote{Em 1887, antes da sua militância como anarquista, Goldman
  conseguiu obter a cidadania estadunidense, ao se casar com o seu
  primeiro marido, Jacob Kerstner. Nos Estados Unidos de então, a
  cidadania da mulher era dependente da cidadania do seu pai ou marido.
  Embora, Goldman tivesse se divorciado de Kerstner cerca de um ano após
  o casamento, a sua cidadania estava garantida. Não obstante, numa
  jogada para lhe retirar o direito de permanecer ou retornar aos EUA
  por parte das autoridades federais que lhe consideravam uma ameaça às
  instituições americanas, a cidadania de Kerstner foi revogada
  postumamente e, com ela, a de Goldman -- que foi deportada para a
  Rússia em 1919. Seis anos após a sua deportação, em 1925, Goldman se
  casou com o cidadão canadense James Colton, de modo a conseguir um
  passaporte e alguma proteção contra a perseguição aos imigrantes.}

\asterisc

Muitas pessoas se mostraram surpresas de que eu, que por tantos anos
critiquei a instituição casamento, tenha, no final, me submetido a ela.

Sem exceções, elas exigem saber se mudei a opinião sustentada no
passado.

Talvez já não soe muito enfática a minha declaração de que agora, como
antes, estou convencida de que a instituição do casamento, do jeito que
é, não acrescenta nada ao que fundamentalmente une homens e mulheres.
Será sempre uma verdade para mim que a união de duas pessoas é um
assunto exclusivamente privado.

\section{Bizarro antes --- Aceito Agora}

Há não muitos anos atrás, esse ponto de vista era considerado chocante e
revolucionário. Atualmente, porém, encontra uma aceitação muito mais
generalizada que que alguns se permitem admitir.

Mesmo os conservadores estão começando a perceber que, muito embora o
casamento possa ser conveniente, não exerce qualquer influência sobre os
impulsos emocionais ou a expressão sexual.

É um fato que o ritual do casamento, propriamente dito, não tem qualquer
relação com a vida e os hábitos dos seres humanos. Na melhor das
hipóteses, a instituição consiste na sanção pública de um acordo privado
entre duas pessoas. Nunca o Estado teve o poder de dar mais do que essa
sanção, pois é sempre o amor que domina esse tipo de relação humana.

Costumo dar risada das antigas referências ao chamado amor livre. Como
se o amor pudesse ser algo que não livre!

Atualmente, milhares de pessoas se submetem à cerimônia de casamento,
não porque acreditem nela, mas porque protege as suas vidas particulares
das bisbilhotices vulgares.

\section{Casamento: a resposta}

E assim, se você me vê casada hoje, você tem, com isso, a resposta -- se
você realmente se importa em ter uma! Um casamento não implica
divorcia-se do próprio ponto de vista.

O maior infrator da sacralidade da privacidade é o Estado. Desde a
reação desencadeada pela Guerra Mundial, o Estado abandonou a maioria
das suas atividades para se dedicar exclusivamente ao exercício de
controle e sufocamento dos indivíduos.

Já não é mais possível ter liberdade de movimento ou de gosto. Inúmeras
restrições cercam o indivíduo da manhã à noite. O que ele come, bebe,
lê, vê, as relações que estabelece, as opiniões que ele ouve e com quem
se reúne -- tudo está sob vigilância constante. O resultado é que o
indivíduo se vê constantemente na necessidade de inventar métodos
através dos quais ele possa escapar dos tentáculos de tais intrusões e
obstruções.

\section{A necessidade é a mãe da estratégia}

Esforços para impedir as manifestações externas sempre têm como efeito a
criação do gênio capaz de inventar modos de escape a partir do interior.
A Proibição serviu apenas para aguçar a sede americana. A censura à
expressão literária teve como única consequência estimular livros cada
vez mais revolucionários. O que também incidentalmente aumentou o
interesse intelectual e a discussão dos trabalhos proibidos.

No que diz respeito a viagens, com a interferência interminável sobre o
deslocamento, imposta pelo passaporte e pelo visto, o Estado obriga as
pessoas a exercitarem toda a sua engenhosidade na arte de atravessar o
muro chinês construído ao redor do mundo durante guerra. Isso não
significa que as pessoas estejam morrendo de amores pelo Estado.

Em certa medida, o mesmo se aplica ao casamento. Ninguém minimamente
inteligente pode ainda acreditar que o casamento é uma obra dos céus, ou
que, na terra, ele pode ser construído por alguém além dos imediatamente
implicados; eles simplesmente atravessam o processo com o mesmo espírito
de alguém que busca tirar um passaporte ou obter um visto -- para
conseguir um espaço para respirar e proteger a privacidade da
personalidade humana.\footnote{Embora no presente texto, Goldman deixe
  suficientemente claro que o seu casamento com Colton se deu por razões
  pragmáticas, muitos chamaram a atenção para a discrepância entre as
  suas declarações anteriores -- como a que o casamento seria uma forma
  ainda mais aguda de prostituição -- e a aqui contida: segundo a qual o
  casamento é um assunto do âmbito exclusivamente privado. Há quem veja
  nessa sua nova ênfase do casamento como assunto privado (não parece
  justo dizer que ela não tivesse sugerido isso desde o seu primeiro
  texto sobre a temática), uma espécie de autojustificativa, como se com
  isso, buscasse evitar cair ela mesma na sua crítica anterior: a do
  casamento por interesse, como única alternativa para a manutenção
  econômica (no seu caso, política) da mulher. Seja como for, em 1927,
  segundo o jornal Toronto Daily Star, Goldman teria dado uma palestra
  em Toronto defendendo o que na época era chamado de ``casamento por
  companheirismo'', uma alternativa bem mais moderada (e um tanto
  reformista) da instituição casamento tal qual havia. Um logo rascunho
  escrito de próprio punho, aproximadamente em 1928, intitulado
  ``Casamento por companheirismo'', confirma o interesse de Goldman na
  temática, pouco depois do seu casamento com Colton.}

\chapter[A luta do feminismo não foi em vão (1926)]{A luta do feminismo não foi em vão: a conclusão de Emma Goldman\footnote[*]{Texto originalmente publicado no
  jornal nova-iorquino \emph{Rochester Times-Union} (1926).}}

Nota do jornal -- A carreira bastante variada de Emma Goldman foi da
denúncia violenta da América e do casamento à sua posição atual, em
Montreal (Canadá), em que na condição de mulher casada, pede para ser
readmitida nos Estados Unidos. Ainda existe outra fase paradoxal da sua
vida -- quando da antiga posição de feminista militante, passou a ser
uma crítica severa do sufrágio feminino. Neste artigo, o terceiro de uma
séria de cinco, são apresentadas as suas opiniões sobre esse assunto de
imenso interesse -- escrito exclusivamente para o \emph{NEA Service} e o
\emph{Times-Union.}

\asterisc

Caso relembremos das profecias radicais das ativistas dos ``Direitos da
Mulher'' sobre os milagres que o feminismo viria a realizar, quando a
mulher finalmente tivesse o direito de sufrágio e igualdade
profissional, é impossível não admitir que os resultados do feminismo
foram tudo menos proporcionais à brava luta empunhada pelas mulheres em
nome da sua emancipação.

Não faz muito tempo, que as grandes líderes feministas nos asseguraram
que o seu credo purificaria a política, aboliria a guerra, eliminaria
todos os males sociais e criaria relações completamente novas entre os
sexos. Hoje em dia, nenhuma feminista inteligente se deixaria levar por
uma conversa tão boba. Finalmente aprenderam que abusos de longíssima
data não podem ser eliminados com o voto de Minerva.

E o que é mais importante: elas aprenderam que a emancipação econômica e
social das mulheres está intimamente ligada à luta pela emancipação
humana em geral -- que a total independência do homem, assim como a da
mulher só ocorrerá com a mudança completa de nossa atual estrutura
social e do valor individual e coletivo.

No entanto, é um fato que a luta heroica empunhada pelas mulheres, já há
tantos anos na América e na Europa, não foi em vão. Se ainda lhe é
negada uma remuneração equivalente pelo trabalho que realiza, a mulher
conseguiu provar que é capaz de realizá-lo muito bem. Não existe
profissão ou ofício, nem mesmo atravessar o Canal da Mancha a
nado,\footnote{Em agosto de 1926, o presente texto foi publicado em
  novembro, Gertrude Ederle se tornou a primeira mulher a atravessar a
  nado o Canal da Mancha.} que seja estranho à mulher.

\section{Destreza mais fina}

Enquanto estive na Alemanha, durante a guerra, pude observar que as
mulheres da indústria metalúrgica possuíam uma destreza maior do que os
homens na fabricação de instrumentos delicados. Que as mulheres durante
o cataclismo mundial não apenas tiveram a capacidade, mas efetivamente
empreenderam as tarefas mais difíceis -- enquanto os homens sangravam
nos campos de batalha --, já não é um aspecto que precise ser
enfatizado, a essa altura do campeonato.

Sim, as mulheres vêm tendo um desempenho muito bom. Ninguém mais ousará
insistir que o seu lugar apropriado é exclusivamente em casa, que deve
desperdiçar sua substância como empregada doméstica ou mercadoria
sexual. Ela rompeu as grades da gaiola dourada e agora se encontra do
lado de fora, no mundo, pronta para assumir a sua parte da
responsabilidade e também para exigir o direito às suas realizações.

As mulheres americanas mais avançadas vêm agindo dessa forma já muito
antes da guerra, mas só agora, depois da guerra, é que as mulheres no
exterior estão conseguindo se recuperar. Existem pouquíssimas
``Gretchens'', dependentes, complacentes, obedientes e submissas na
Alemanha de hoje. As mulheres na Inglaterra tampouco estão se valendo de
sigilo e subterfúgios para lidar com a vida e as condições sociais.
Aberta e francamente, elas declaram o seu direito a qualquer
conhecimento e experiência que exista no mundo.

E mesmo na França, as mulheres, além do direito de amar, sobre o qual
alegam maestria, estão começando a perceber que a vida é mais do que
namoricos fúteis, que há problemas sociais muitos graves e que,
portanto, exigem a atenção das mulheres, assim como a dos homens. Em
outras palavras, em todo o mundo, as mulheres se tornaram profundamente
conscientes da necessidade de fazer a sua parte na grande luta mundial.

\section{Vital em todos os domínios}

A mulher atual é, possivelmente, a força mais vital em todos os domínios
do pensamento e empreendimentos humanos.

Se tal condição é efeito da perda de vitalidade de muitos homens, por
conta dos horrores da guerra, é algo que não posso saber. O que eu sei é
apenas que a maioria dos homens de classe média da Europa perdeu o
controle sobre a sua vida. Eles não têm mais fé ou idealismo. Para usar
uma expressão ``freudiana'', a maioria dos homens de hoje parece sofrer
de um complexo de inferioridade. Ou se trataria de orgulho ferido por
não mais poder bancar o bravo cavalheiro que protegia a mulher de viver
tão perigosamente quanto eles?

De uma forma ou de outra, o ponto é que, em geral, os homens parecem
estar perdidos, ``desempregados'', de certa forma. Eles não sabem o que
fazer com eles mesmos na presença daquelas que antes eram consideradas
inferiores.

\section{Viva, Ávida e Ativa}

Não é isso o que ocorre com as mulheres que conheci na Europa. Elas me
impressionaram pela completa transformação das suas qualidades físicas,
mentais, espirituais e emocionais -- um tipo novo e viril de
feminilidade, muito mais vivo, ávido, ativo e livre do que o dos homens.

Muitos fatores contribuíram para criar o tipo moderno de mulher, sendo o
fator mais vital a solidariedade das mulheres entre si. A necessidade as
ensinou, já na fase inicial da sua luta, que o escravo nunca foi liberto
pelas mãos do seu senhor e que a sua emancipação só poderia ser
conquistada pelo espírito de solidariedade entre os seus companheiros
escravos.

Do mesmo modo, a solidariedade do sexo entre as mulheres tem propiciado,
creio eu, um tremendo ímpeto e encorajamento na luta pela autoafirmação
e pelo direito do seu lugar no mundo -- o direito de ser ela mesma.

\chapter[O elemento sexual da vida (1935)]{O elemento sexual da vida\footnote[*]{Rascunho inacabado, do qual muitas
  páginas estão faltando, a datação de 1935 é aproximada. Ao que parece foi
  concebido como parte de um texto de abordagem mais ampla, a tirar pelo
  título: ``Sexualidade, maternidade e controle de natalidade''. Para
  além da ``seção'' aqui traduzida, bastante fragmentada, não restaram
  mais do que duas ou três páginas do documento original. As
  pouquíssimas passagens do rascunho excluídas da presente tradução ou
  foram riscadas pela própria Goldman (há uma série de anotações à mão
  no texto, nem sempre legíveis), ou a redação não revisada inviabilizou
  a inteligibilidade do texto.}}



A verdade virá à tona algum dia. Mas, em geral, é a mentira que
persiste. A verdade é nua, direta. Não se apresenta com sorrisos
maliciosos, subterfúgios ou acordos. Veja as mentiras adornadas com seda
e joias. Elas são insinuantes, lisonjeiras e enganosas. Muitos,
incontáveis, ficam deslumbrados com a pompa e autoimportância com que se
transvestem as mentiras e, assim, seguem-nas, alegremente, sem perceber
a face irônica sob a máscara cheia de enfeites.

Não é motivo de surpresa, então, que a força mais elementar da vida
humana, o sexo, seja ainda hoje degradado e renegado.

As duas instituições que, há séculos, tentam subjugar o sexo, arrancá-lo
através dos métodos mais perversos, são a igreja e a moralidade -- as
caluniadoras de tudo o que é excelente e saudável na vida. Não obstante,
quanto mais a igreja e a moral tentaram subjugar o sexo e extingui-lo
das necessidades do homem, mais intensa e devastadoramente o elemento
sexual se afirmou.

É relativamente recente que a verdade sobre o sexo rompeu a rede de
falsidades, ilusões e armadilhas que há tanto tempo assombram a mente do
homem.

Os grandes psicólogos do sexo, Havelock Ellis, Kraft Ebbing, Edward
Carpenter e Freud estavam sozinhos na sua batalha heroica contra as
mentiras sobre o sexo. Atualmente, o número pode ser contado aos
milhares.\footnote{Boa parte desses dois períodos, com os quais o
  parágrafo é iniciado, foram riscados e corrigidos à mão por Goldman.
  Como, porém, a sua caligrafia nessa passagem, não está com boa
  legibilidade, optou-se por manter tal qual datilografado.} E todos
comprovam, para além de qualquer sombra de dúvida, que as fobias, os
terrores, as neuroses e demais transtornos mentais têm a mesma origem
sexual. Mais do que isso: toda a nossa vida, desde as nossas ações mais
simples, depende completamente dos nossos impulsos sexuais. A
dissimulação imposta pela longa repressão moral nos impediu de
reconhecer essa verdade fundamental. Não obstante, na medida em que é
impossível resistir impunemente a uma força natural tão poderosa, nós
nos tornamos todos mais ou menos insanos: alguns completamente, outros
expressaram suas aspirações sexuais através de fobias passageiras ou
sonhos monstruosos; outros, ainda, na tentativa de se libertar dessas
aspirações, passaram quase toda a sua existência como se sonhando
acordados, indiferentes à vida cotidiana. Por fim, aqueles que eram mais
fortes satisfizeram suas propensões valendo-se de uma hipocrisia
calculista que lhes permitia se ocultar dos olhos da multidão.

O crescente acúmulo de pesquisas de abordagem freudiana, as
consequências, cada vez maiores, da aplicação das doutrinas originais,
fornecem explicações para muitos dos problemas que costumavam nos
confundir e, frequentemente, possibilita-nos encontrar provas materiais
da teoria da unidade de energia.\footnote{Possivelmente, Goldman esteja
  se referindo aqui à teoria da libido.} Efetivamente, vemos que o
sonho, o sono, o estado normal de consciência e as psicoses formam uma
parte integrante da personalidade humana, que reage às influências
externas via uma ou outra dessas manifestações, de acordo com a fase da
evolução dessa força no momento. Observamos que dentre essas forças que
atuam sobre nós, a mais importante, se não a única, a que quase sempre
encontramos sob os mais variados disfarces, é o impulso sexual.

Precisamos admitir que não é uma possibilidade negar a realidade, mesmo
que não gostemos dela, e, nesse sentido, é indispensável que
reconheçamos no tão difamado impulso sexual a força psicológica motriz
da humanidade.

\asterisc

O falecido Prof. Dorsey, no seu livro ``Por que nos comportamos como
seres humanos'', chamou a atenção para o que a psicologia diagnosticou
como ``complexo de impureza'' e nos mostrou o que está por trás desse
puritanismo descarado que anuncia a ``pureza'' desse ou daquela. Ele
também demonstra que a pureza do ignorante, quando adquirida pelo preço
da repressão da sua curiosidade natural, não é segura e nem sã. De outro
lado, continua o professor Dorsey, o estudo da biologia está começando a
derrubar esse ``complexo de impureza'' e a doutrina profana e
antinatural, iniciada pelos primeiros monges cristãos, de que o impulso
sexual é o sinal de degradação do homem e a fonte da sua energia mais
diabólica. A natureza sempre sabe mais.

O sexo é uma função biológica primária a toda forma de vida superior.
Suas características e qualidades vêm de uma linhagem antiga. Seu
impulso é tão real, quanto é a força que faz as marés baixarem e
subirem. Possui estrutura e comportamento profundamente suscetíveis. É o
elemento fundamental de toda a vida superior, seus caracteres externos
são um engenhoso estratagema publicitário da natureza, através do qual
ela vende seus produtos e, assim, assegura sua família.

Devemos ao sexo mais do que à poesia; devemos o canto dos pássaros, toda
música cantada e a própria voz, devemos a plumagem que atinge a glória
suprema com as aves-do-paraíso, devemos a juba do leão e todas as formas
superiores da vida tanto no mundo vegetal, quanto animal. Está tecido em
todas as estruturas da vida humana e inserido em todos os costumes. Na
conta do sexo, devemos debitar muitas tolices e estupidezes e algumas
das mais insensatas injustiças que tornam essa coisa que chamamos de
cultura humana o mosaico incrível e variado que é.

É verdade que somos mais esclarecidos do que éramos, mas ainda não
atingimos o estágio em que uma simples menção ao sexo não provoque
respostas de reprovação ou insulto. Não são poucas as pessoas que
defendem que o ``conhecimento sexual'' é algo questionável, ou, na
melhor das hipóteses, algo que deve ser mantido em sigilo no interior
dos ``livros de medicina''; ou, ainda, há quem considere a posse desse
conhecimento uma banalidade condizente com os irremediavelmente
desavergonhados. Como resultado, a literatura pseudocientífica acerca do
``sexo'' é derramada sobre as nossas emoções, deixando os fatos
intocados, quando não os apresenta sob disfarces. Muito do que é
produzido não possui fundamentação biológica ou qualquer relação com as
leis da vida que governam o homem não menos do que a qualquer outro ser
vivo. É o medo (às vezes chamado de ``reverência'') que nos faz ``deixar
o sexo em paz''. São o recato dissimulado e a vergonha estúpida,
disfarçados sob o nome de ``decência'', que obrigam os museus a vestir
faunos de mármore e a cobrir júpiteres e cupidos de bronze com folhas de
gesso, não raro tortas ou mal cortadas dos lados. Há considerável
infantilidade a envolver a questão do sexo.

O homem é ``superior'', os animais são ``inferiores'' -- sem intelecto
e, é claro, não podem ter ``almas''. Nós temos. O nosso pai é um pai
``divino''. Nossos corpos são ``sagrados''. Consequentemente, a arte, da
escultura de Fídias ao poema sofomórico {[}\emph{sophomoric poem}{]},
tende à glorificação cada vez maior do ser humano: homens e mulheres
mais próximos a deuses e deusas; deuses e deusas como homens e mulheres
glorificados.

E assim aconteceu de que tudo aquilo de mais comum na natureza, além de
vivo, se tornasse dotado da santidade do céu.

\asterisc

A fisiologia e a zoologia nos mostram que a raiz da beleza e do sentido
estético reside no elemento sexual -- que não pode e nem deve ser
ignorado. A vida sexual, da forma mais primitiva em diante, busca a
união, a coesão. Em toda parte, há signos de atração: nas plantas,
tem-se as cores vibrantes e as formas das flores; nos pássaros, a doçura
vitoriosa do seu canto. O pássaro pavilhão da Austrália decora a sua
``dança'' com flores e plumas; o pavão fascina sua parceira; os alces
gemem um para o outro nas florestas; os vagalumes acendem suas lanternas
à noite; o ar está cheio de odores místicos flutuando. Todas as
faculdades e formas da natureza estão dispostas de modo a contribuir com
a expressão dessa grande necessidade de união que explode no mundo
animal. Tudo se transforma em indício, símbolo, lembrança, mensagem,
chamado. Essa faculdade é social. É o começo do panorama da arte.

Penso que não cometeremos um erro ao afirmar que a aceitação livre e
saudável do corpo humano, com todas as suas faculdades, será a chave
mestra da arte do futuro.

\asterisc

Todos os escritores modernos que tratam da questão do sexo provaram que
a velha noção de que o sexo se inicia na puberdade, é falsa. O impulso
sexual, como todos os nossos impulsos, começa no nascimento e termina
com a morte. Diferentemente dos outros animais, porém, ao homem não é
permitido interpretar que os seus apetites sexuais exigem a ação, ao
menos não antes que ele atinja um número respeitável de anos (uma causa
verdadeiramente excelente para os distúrbios nervosos dos quais os
demais animais, fiéis aos retornos periódicos da sua excitação sexual,
estão absolutamente livres). Ao que parece, há um despertar da
sexualidade desde os primeiros dias de vida; para Freud e sua escola, as
ideias sexuais ganham existência no momento em que é dado à criança o
peito.

Aparentemente, não é só a sensação de satisfação do apetite, o que
impele o recém-nascido a agarrar-se ao seio da sua mãe ou da ama de
leite. Nessa pequena bola disforme, na qual a gama de sensações ainda
não emergiu na sua totalidade, uma imaginação especial e antecipatória
já complica o prazer de ter o apetite satisfeito por meio de um gozo
mais especializado.

E, passará um bom tempo dos primeiros anos em que a imaginação da
criança está voltada para imagens sexuais. Isolada do mundo e sem saber
nada além da vida familiar, embora já possuidora de desejos que a
natureza mesma não lhe permitirá satisfazer até um futuro distante, a
criança volta a sua imaginação sexual para os objetos mais próximos, os
únicos que pode ter em vista. Há, na primeira infância, uma crise
imaginativa que, frequentemente, exerce um efeito distante sobre a vida
sentimental do adulto.

Os psicólogos nos pedem para observar que o desenvolvimento sexual da
criança pode ocorrer em três etapas:

\begin{enumerate}
\item Surgimento do primeiro fenômeno no momento do nascimento;

\item Crise na primeira infância, de quatro a sete anos, com a inversão de
sentimentos normais, incesto, etc;

\item Crise da puberdade, entre dez e quatorze anos, seguida pelo retorno
das sensações às vias normais.
\end{enumerate}

Essas manifestações psicológicas estão, sem dúvida, intimamente ligadas
às modificações dos órgãos de secreção internos, e, em particular, com o
desenvolvimento sexual, que é apenas uma das vertentes dessas atividades
glandulares.

Antigamente, todas essas manifestações eram recebidas com profunda
ignorância e estupidez. Repetidas vezes, as mães bateram e ainda batem
nos seus desafortunados filhos, quando veem a menor manifestação do
sexo. Ou elas os aterrorizam, colocando o estigma da depravação nas suas
crianças. Eu poderia citar centenas de exemplos retirados da minha
experiência de enfermeira formalmente treinada.

A sociedade exige que jovens adultos, homens e mulheres (especialmente
as mulheres), reprimam o impulso sexual por vários anos -- geralmente
por toda a sua vida. A contenção de um desejo tão instintivo não pode
ser bem-sucedida numa pessoa normal sem consequências diretas na sua
saúde -- todo tipo de distúrbio mental e físico pode resultar daí; e não
é raro que o impulso, forte demais para ser contido, encontre a saída
por alguma via infantil e perversa.

\asterisc

Ou conforme foi tão habilmente descrito na obra de Jacques Fisher,
\emph{Amor e Moralidade}: o desejo sexual frequentemente estoura, como
se fosse um trovão, nas vidas mais ordenadas e respeitosas, que, à
primeira vista, pareceriam as mais imunes. Algumas vezes, sob a
influência de uma causa fisiológica específica (influência da menopausa
ou de alguma alteração nas secreções de um órgão prejudicado pela idade)
e outras de uma causa desconhecida. As regras morais que até aquele
momento governavam toda a vida sexual desaparecem ou então sofrem uma
revolução violenta.

Infelizmente pouquíssimos entendem tal mudança. A eles não foi permitido
ver além dos atos propriamente ditos. Eles não se dão o trabalho de
descobrir os motivos que conduziram {[}até ali{]}.\footnote{Parágrafo
  riscado no manuscrito datilografado.}

\asterisc

Max Stirner -- \emph{O único e sua propriedade} -- oferece um retrato
maravilhoso da sexualidade faminta.

``Para onde quer que se olhe, encontraremos vítimas da autonegação. Ali,
na minha frente, está uma garota que talvez já há dez anos vem fazendo
sacrifícios cruéis à sua alma. Sobre o seu corpo cheio de vigor pende
uma cabeça cansada de morte, e as maçãs pálidas do seu rosto denunciam a
sangria lenta da sua juventude. Pobre criança, quantas vezes as paixões
bateram na porta do seu coração e os intensos poderes da juventude
reclamaram os seus direitos? Quando a tua cabeça se revirava no
travesseiro macio, como a natureza desperta tremia pelos teus membros,
como o sangue inchava as tuas veias e fantasias ardentes derramavam o
brilho da volúpia nos seus olhos! Mas aí aparecia o fantasma da alma e
da sua bem-aventurança eterna. Você ficava aterrorizada, suas mãos se
recolhiam, seus olhos atormentados se voltavam para o alto, você --
rezava. As tempestades da natureza foram silenciadas, a calmaria
deslizou sobre o oceano dos seus desejos. Lentamente, as pálpebras
cansadas caíram por sobre a vida que se apagou debaixo delas, a tensão
foi deixando, imperceptivelmente, os seus membros bem torneados, as
ondas turbulentas secaram no seu coração, as mãos unidas em posição de
oração quedaram, como um peso sem força, sobre um peito que não opunha
resistência, um último suspiro ``Ó, querido'', saiu como lamento, e -- a
alma estava enfim em paz . Adormecias, e despertavas pelas manhãs para
novas batalhas e novas -- orações. Agora, o hábito da renúncia gela o
calor do teu desejo e as rosas da tua juventude estão se tornando
pálidas na -- anemia da sua beatitude. A alma está salva, que o corpo
pereça. Ó, Lais, ó Ninon, quão bem vocês fizeram em desdenhar dessa
pálida virtude! Uma única \emph{grisette} livre contra mil virgens cuja
virtude tornou murchas!''

Não há nada mais patético e cruel do que a vida de uma solteirona ou de
uma celibatária que murchou de tanta virtude. Com seus pensamentos e
sentimentos reprimidos, elas se tornam mal-humoradas, invejosas e
ciumentas. Ressentem-se da despreocupação e alegria dos jovens. Mas será
que podemos as culpar por isso? Certamente, não. Elas são vítimas de uma
moralidade sexual ridícula.

É certo que a repressão sexual levada a cabo por um longo período de
anos mesmo que não seja suficiente para causar uma desordem mental
completa, cria uma série de peculiaridades antissociais e perversões.

Embora alguns moralistas enfatizem que o sexo é uma força dominante em
todas as criaturas vivas -- seja um ser humano, animal ou planta --,
eles têm para si a ideia de que o sexo é mais forte no macho do que na
fêmea. Também insistem em atribuir ao sexo apenas uma única motivação --
a reprodução. Essas ideias foram exploradas pelos psicólogos sexuais
modernos que, por sua vez, insistem que a diferença em intensidade é
normalmente exagerada ao manter longe das garotas o conhecimento de
questões sexuais, ou ao sugerir que há algo mau nesses assuntos. O sexo
é uma desgraça para as boas meninas.

O célebre ginecologista, Kisch, nos diz:

``O impulso sexual é extremamente poderoso em certos períodos da vida, é
uma força elementar que domina de modo tão absoluto a totalidade do
organismo de uma mulher que não deixa na sua mente lugar algum para
pensamentos acerca da reprodução; pelo contrário, ela deseja
intensamente o ato sexual mesmo quando está com muito medo de engravidar
ou quando não há chance alguma de a gravidez ocorrer.''\footnote{Enoch
  Heinrich Kisch. \emph{The Sexual Life of Woman in Its Physiological
  and Hygienic Aspect}. (1910)}

Já o Dr. Bloch escreve:

``Eu mesmo entrevistei um número significativo de mulheres cultas acerca
desse assunto. Sem exceção, elas declararam que a teoria de que a mulher
possui uma sensibilidade sexual menor é equivocada; muitas eram
inclusive da opinião que a sensibilidade sexual era maior e mais
duradoura na mulher do que no homem.''

Dr. Bloch também insiste que: ``A maioria dos casos de frigidez sexual
nas mulheres é, na verdade, meramente aparente: seja porque, sob o véu
prescrito pela moralidade convencional, por detrás de uma frieza
aparente, há uma sexualidade ardente oculta, ou ainda porque o homem com
quem ela teve relação sexual não foi bem-sucedido em despertar a sua
sensibilidade erótica, tão complicada e difícil de ser desperta. Não
obstante, uma vez que ele tenha sucesso em despertá-la, a
insensibilidade sexual, na maioria dos casos, simplesmente,
desaparecerá.''\footnote{Iwan Bloch. \emph{The Sexual Life of Our Time
  in Its Relations to Modern Civilization.}(1908)}

Na verdade, muitas esposas não se atrevem a dar o seu máximo por medo de
que seus maridos as julguem muito agressivas, desprovidas do tipo
correto de feminilidade. A maioria dos homens são criados para acreditar
que as mulheres devem ser possuídas e que jamais devem entregar a si
mesmas ao amor e à paixão com felicidade e alegria.

Isso também impede que homens de uma espécie mais sensível possam se
entregar livremente -- eles temem ultrajar e chocar a sensibilidade e
inocência das suas esposas. Vocês se surpreenderiam com a frequência que
esposas se sentem chocadas e ultrajadas.

Laura, esposa do capitão do drama \emph{O pai} de Strindberg, nos
oferece uma chave para compreender os pensamentos e sentimentos de um
grande número de mulheres. Ela diz para o seu marido: ``Quando,
primeiramente, eu entrei na sua vida, eu era como uma segunda mãe para
você. Eu te amava como meu filho. Mas quando a natureza dos seus
sentimentos mudavam e você surgia como o meu amante, eu me envergonhava
e seu abraço era um prazer logo seguido por uma consciência cheia de
remorsos como se meu sangue estivesse envergonhado.'' Essa é a tragédia
de muitas mulheres.

\asterisc

O que é comumente chamado de incompatibilidade de temperamentos é quase
sempre o resultado direto da ausência de harmonia sexual -- a
insatisfação e os atritos que surgem quando a natureza química do sexo
entre marido e mulher falha em os unir harmoniosamente.

A frigidez, em algumas mulheres, é em grande medida um efeito do
sufocamento imposto pelo tabu sexual. Essas mulheres não podem, mesmo
ainda que tentem desesperadamente, responder ao desejo sexual do homem.
Na verdade, só o pensamento da relação sexual é, para tais mulheres, uma
tortura. Mesmo que falte ao homem a sutileza de modo que ele imponha
suas necessidades à esposa, não encontrará de nenhuma satisfação. No
final, ele buscará gratificação em outro lugar. Há uma porcentagem
considerável de homens casados entre a clientela da prostituição. O sexo
é mais poderoso do que todas as decisões. O homem se tornará indiferente
e, por fim, desejará o divórcio ou se ele puder pagar poderá manter uma
amante.

\asterisc

O Sr. V. F. Calverton no seu livro \emph{A expressão sexual na
literatura} {[}\emph{Sex expression in literature}{]} mostra que ``o
melhor método para combater a mitologia popular e antiga de que nossas
atitudes e instituições sexuais são originalmente reveladas ao homem por
deus, de modo que devem permanecer fixas, imutáveis e eternas, é, sem
sombra de dúvida, a execução de um estudo realista sobre a história da
moralidade sexual em todas as diferentes eras culturais.''

Ele prova que remanescentes do culto ao sexo eram comuns até o fim do
século XVIII -- as festas fálicas de Dionísio foram mantidas em muitas
províncias da Itália. Para citar as palavras de W. J. Lawrence, ``quando
a cidade antiga do reino de Nápoles, Isérnia, foi destruída por um
terremoto em 1805, os crentes viram na calamidade um castigo para os
costumes pagãos conservados inescrupulosamente no período do ano do
Festival Báquico.''\footnote{W. J. Lawrence. ``The Pallus on the Early
  English Stage''. In: \emph{Psyche and Eros. An International Bi-Montly
  Journal of Psychanalysis, Psychology and Therapeutic Psychognosis,
  vol. II n. 3.}}

... Depois da conversão ao cristianismo, as raças germânicas mantiveram
os seus demônios fálicos da vegetação. Nos tempos antigos, as mulheres
de Israel faziam falos de ouro e prata para serem usados como talismãs
contra a esterilidade; também havia os doces em forma de falo que foram
comidos ao longo de todo o período das trevas com o mesmo propósito em
vista.\footnote{Neste parágrafo, Goldman faz um pequeno resumo de
  algumas das ideias encontradas num outro artigo publicado no mesmo
  volume que o de Lawrence citado no parágrafo anterior, sendo este:
  ``Sex Fertility and Worship'', cujo autor, anônimo, identifica-se como
  R.}

Tanto a religião judaica, quanto a cristã impuseram a noção de que o
sexo só é admissível para o ato de procriação. Para que pudessem ser
frutíferos, se multiplicassem e reabastecessem a terra, Jeová permitiu
aos judeus ter muitas esposas. A religião cristã, embora não permita
muitas esposas e tenha aceito a gratificação sexual apenas com a
finalidade de trazer crianças ao mundo, está perfeitamente consciente do
fato de que seja dentro ou fora do casamento, o sexo continua a existir
sem muita consideração pela procriação. Mas enquanto muitos religiosos e
puritanos predem-se ao fetiche de que o sexo não deve ser ``indulgente''
-- conforme eles nomeiam essa expressão perfeitamente natural --,
psicólogos e biólogos modernos já rasgaram os véus de todo absurdo que
envolve o sexo. Bibliotecas inteiras foram preenchidas com trabalhos que
tratam do assunto com entendimento e profundidade e demonstraram que
houve uma percepção inadequada, de um lado, da energia gigantesca que
está por detrás do instinto sexual e, de outro, dos dispositivos
biológicos para a liberação dessa energia através de canais que não são
especificamente sexuais. Provavelmente, é uma verdade absoluta dizer que
o instinto sexual é o instinto criativo, como é igualmente verdadeiro
que qualquer escoadouro que ofereça satisfação emocional oriunda do
exercício da criatividade tem a capacidade de neutralizar as
necessidades dos desejos sexuais.

O significado desse ponto de vista deve ser plenamente concretizado
pelos pais e seus filhos, pois as crianças são seus reflexos não apenas
por conta do fator hereditário, mas também pela assimilação das suas
ideias e ideais. Ambos devem entender e aprender a utilizar aquilo que é
essencial para desenvolver e satisfazer os desejos emocionais oriundos
do instinto sexual. É nessa relação que a liberação da energia, por meio
de atividade recreativas e de desenvolvimento de habilidades, dá lugar
aos resultados mais satisfatórios.

O espírito criativo não é, porém, um antídoto para o instinto sexual,
antes deve ser considerado uma uma parte da sua expressão mais vigorosa.
Age de forma a conservar, mas também utiliza o instinto para formas de
satisfação que não são tem por objetivo meramente a proteção, já que o
conduzem a um desenvolvimento ainda maior. Seu alargamento e
aprofundamento se imprimem sobre o o caráter inato da pessoa e seu poder
de autodirecionamento e controle. A liberação não-sexual de energia,
algumas vezes, é suficiente para compensar as necessidades fundamentais
do desejo sexual e, ainda, para transmutá-las em autossatisfação e
formas úteis de expressão.

O homem que expressou esse pensamento do modo mais profundo e poético
foi Friedrich Nietzsche. Em ``Moral como antinatureza'' ,\footnote{Título
  de um dos capítulos da obra \emph{Crepúsculo dos ídolos.}} ele
escreve:

``Todas as paixões têm uma época em que são unicamente destrutivas,
quando, com o peso da sua estupidez, puxam suas vítimas para baixo; e
têm um período posterior, muito posterior, em que se casam com o
espírito, em que são ``espiritualizadas''. Outrora, as pessoas
declaravam guerra contra as próprias paixões, em virtude da sua
estupidez, conspiravam para a sua aniquilação. A formulação mais notável
desta concepção encontra-se no Novo Testamento, no Sermão da Montanha,
onde, diga-se de passagem, as coisas não são, de modo algum,
consideradas de um ponto de vista elevado. Por exemplo, com relação à
sexualidade, ali está posto de: `Se o teu olho te escandalizar,
arranca-o fora'.''

Felizmente, nenhum cristão age de acordo com esse preceito. Aniquilar as
paixões e desejos unicamente para prevenir a estupidez e resultados
desagradáveis é, para nós, no presente, simplesmente uma forma ainda
mais aguda de estupidez. A igreja luta contra as paixões, através da
extirpação em todos os sentidos; essa é a sua prática, a sua ``cura'', a
sua castração. Ela nunca questiona: ``Como espiritualizar, embelezar e
deificar um desejo?'' -- em todos as eras, ela sempre colocou a ênfase
na disciplina para a exterminação (da sensualidade, do orgulho e
ambição). No entanto, atacar as paixões na raiz significa atacar a
própria vida na raiz: a práxis da igreja é inimiga da vida...

Esse truísmo já é reconhecido pelo pensamento, mas um elemento daquela
etapa ainda permaneceu, a saber: que o sexo, embora um fator
extremamente potente, deve ser reprimido por aqueles que são solteiros.
Felizmente, esse preconceito também está sendo demolido. Em um simpósio
recente sobre ``o elemento sexual na vida do adulto solteiro'' foram
apresentados fatos claros e surpreendentes sobre o assunto. Para citar
apenas alguns dos colocados pelo Dr. Ira Wile na sua contribuição:

``Os solteiros possuem os mesmos potenciais das pessoas casadas com as
quais estão necessariamente relacionados nos interesses e na
participação em todos os aspectos da vida e, portanto, afetam o
bem-estar do grupo de casados. A sua vida sexual, de várias formas, é
fundamental para o bem-estar social, assim como é significativa para o
seu crescimento pessoal e desenvolvimento.

``Estudar essa parte da sexualidade é explorar uma das áreas não
mapeadas da nossa civilização na qual o resíduo da ignorância bárbara e
do medo ainda retardam o progresso. Acima das suas profundezas, aparece
uma névoa miasmática e pouco convidativa de dúvida e incerteza. Até
agora, poucos exploraram de modo sistemático esse estado dos solteiros.
Consequentemente, o esforço de fazer uma pesquisa nesse campo, abordando
os seus vários ângulos, representa um esforço de mapear fatos, rastrear
influências casuais, determinar a validade e o valor das nossas visões
atuais e estabelecer dados e hipóteses que possam estar a serviço de
interpretar a nossa era. {[}...{]}

Há mais de dois mil anos a recompensa oficial à urgência de procriação
tem sido relegada ao estado de casado socialmente reconhecido. O
pressuposto da sociedade é que a proibição autorizada ao desejo por um
parceiro biológico deve ser capaz de adiar a sua efetivação até que o
rito civil ou religioso exigido sancione o desejo. Se a sociedade espera
que regulamentações externas subjuguem necessidades fisiológicas e
impulsos psicológicos é o caso de perguntar se essas regulamentações têm
sido efetivas. Aqueles que não estão casados vivem em castidade e
celibato? Se não, em quais aspectos e como o seu comportamento sexual
difere daquele exibido pelos que lhes são semelhantes em todo o resto,
exceto por viver em casamento?''\footnote{Ira S. Wile. \emph{The sex
  life of the unmarried adult: an inquiry into and an interpretation of
  current sex practices.} (1934)}

A atividade sexual não é um ato isolado; é uma experiência generalizada
que motiva e afeta a personalidade. De reações completamente pessoais
emergem ideias de romance e beleza, exaltação e tranquilidade, devoção e
idealização servil; ou, o sentido de sacrifício, confusão, humilhação,
vergonha, ansiedade, o desejo por autopunição, ou a autoaceitação da
fraqueza, fracasso e inadequação. O conceito de ``apaixonar-se'' pode
ser aplicado a si mesmo, a membros do mesmo sexo ou do sexo oposto. Há
normalmente uma segunda pessoa que facilita o crescimento, seja dos
impulsos sexuais, ou da própria vida pessoal como um todo. Na mesma
medida em que isso é verdade para aqueles que são solteiros, é para
aqueles que são casados, embora a sociedade torne essa harmonia um
desafio, uma propriedade, uma questão.

Havelock Elis ao se referir à reprodução, salienta que: ``As funções do
sexo no que diz respeito aos aspectos psíquicos e emocionais estão muito
além de qualquer ato de procriação, podem inclusive exclui-lo
completamente; e quando a nossa preocupação está relacionada ao
bem-estar do ser humano enquanto indivíduo, devemos ampliar a nossa
perspectiva e aprofundar a nossa visão.''\footnote{Havelock Ellis,
  \emph{Studies in the Psychology of Sex: Sex in Relation to Society},
  vol. 3.}

Para interpretar a vida sexual dos solteiros, é preciso reconhecer que
há duas funções do sexo:

Uma, a biológica, tem a procriação como fim -- envolve um processo cujo
interesse na perpetuação da espécie é muito mais emocional do que
intelectual.

A outra função consiste na promoção do desenvolvimento social através
das relações humanas. Isso envolve o sexo e atividade erótica, com ou
sem, o fim da procriação.

Há duas bases para a energia do impulso sexual -- uma é consciente,
direcionada, guiada, submetida ao controle ético; a outra é
inconsciente, instintiva, impulsiva, responde a estímulos e não está
sujeita à razão.

\asterisc

Numerosos estudos revelaram que tanto os solteiros, quanto os casados
têm experiências sexuais variadas. Dr. Klatt descobriu que 18\% das
mulheres tiveram alguma forma ativa de experiência sexual com menos de
18 anos. E que 60,6\% entre 18 e 22 anos. Entre o grupo seleto do Dr.
Hamilton, 59\% dos homens e 47\% das mulheres fizeram sexo antes do
casamento, e 20\% dos homens e 16\% das mulheres antes dos 21 anos de
idade. Katherine B. Davis, ao estudar a vida sexual de 1200 mulheres
universitárias solteiras, relatou que 61\% admitiram se masturbar, sendo
que 57\% antes dos 15 anos de idade.

Sentimentos e desejos sexuais são experimentados com periodicidade por
70 por cento do grupo reportado. Os solteiros, não menos do que os
casados, relataram várias formas do fenômeno sexual. É interessante
notar que relações homossexuais foram reportadas mais frequentemente
pelo grupo de solteiros. Isso é significativo até porque, socialmente, a
homossexualidade é considerada mais perigosa entre homens do que entre
mulheres. Boa parte das experiências homossexuais entre os homens finda
quando ganham confiança na sua potência sexual, como confirmado na
relação sexual que é seguida pelo casamento.

O lesbianismo é comum e a sua prevalência é importante, ainda que nem
sempre implique em permanecer solteiro. O grupo de mulheres solteiras
lamentou mais a falta de filhos do que de maridos. Aproximadamente,
38,7\% acredita que a relação sexual é necessária para a saúde mental e
física. Aproximadamente, 20\% se pronunciaram a favor da relação sexual
antes do casamento para homens e para mulheres.

Dr. Wile pergunta:

``Os desejos sexuais de adultos solteiros das graduações e
pós-graduações das universidades dotados de menos intensos são, em
geral, menos intensos do que entre os estudantes do ensino médio? Os
jovens mais maduros perdem os seus desejos instintivos primários na
quando se tornam adultos? A educação superior é adquirida à custa da
repressão sexual? Isto é absurdo e contrário a todos os fatos. A
atividade sexual se dá de muitas maneiras -- por meio de jogos eróticos
ativos, praticas autoeróticas, práticas homossexuais ou heterossexuais,
ou via a sublimação estética ou vocacional. A sua natureza e intensidade
estão sujeitas às escolhas, julgamentos, padrões e ideais pessoais que
regulam, mas não aniquilam -- temporalmente proibidos ao invés de
permanentemente inibidos.

Se alguém assumir que a vida sexual é e deve ser limitada pela
instituição social do casamento, então é o caso de perguntar o que fazem
os solteiros ou o que eles devem fazer? A questão não é sobre o que eles
podem fazer, posto isso ser idêntico em ambos os grupos. A existência
real e legal do pacto antenupcial atesta o fato de que há efetivamente
uma atividade sexual bastante abrangente na vida de
solteiro.''\footnote{Ira S. Wile. \emph{The sex life of the unmarried
  adult.}}

\asterisc

A dificuldade econômica não é capaz de erradicar desejos heterossexuais,
como tampouco a promulgação de uma lei punitiva pode destruir impulsos
homossexuais. Na verdade, a sociedade reconhece as demandas sexuais, se
não as necessidades sexuais mesmas, da parte do grupo dos solteiros,
dada a sua atitude em relação à regulamentação da prostituição, da
perseguição aos homossexuais e das casas de strip-tease.

Em outras palavras, ou os solteiros são levados a métodos artificiais
para a expressão sexual ou eles se entregarão ao desejo que domina as
relações naturais. Na realidade, eles têm se entregado às relações
sexuais ainda que sob os tabus que lhes são impostos há tanto tempo, e
têm se entregue especialmente desde que começaram a ter conhecimento dos
métodos contraceptivos, capazes de prevenir que se traga crianças
indesejadas a esse nosso mundo maravilhoso -- um assunto sobre o qual eu
pretendo falar antes de ir embora dessa cidade.

Adultos solteiros estão lidando com o sexo mais como um fato do que como
teoria. Eles estão aceitando a sua compleição sexual abertamente como um
instrumento de crescimento pessoal e emocional, que dá lugar à
estabilidade social, ao invés de assumirem a hipocrisia de que o sexo é
uma função projetada em algum plano divino com o fim da procriação de
seres puros cujo sentido da vida é o de morrer em castidade para, assim,
atingir a felicidade em um mundo vindouro.

Eles valorizam o sexo como a fonte da vida, pois acreditam que uma vida
sem sexo é uma palhaçada após a maturidade biológica, porque é contrária
à natureza, uma vez que o sexo é também expressão de fatores
psicológicos cuja função é socializar os indivíduos e os conduzir a
abandonar a vida de solteiro em nome de um casamento que esteja em
harmonia com as leis e costumes da sua época.

Todos os jogos eróticos estão geneticamente ligados ao namoro, mas o
namoro está sujeito ao controle social. Por conta disso, a vida sexual
de um jovem adulto solteiro e normal, seja lá o que isso queira dizer, é
a preparação para a perfeição da relação sexual, para a promoção da
felicidade pessoal, e para a adaptação a alguma forma de casamento seja
essa livre e não convencional ou de acordo com os ritos civis ou
religiosos.

É por concordo tão completamente com esse ponto de vista e por saber dos
resultados trágicos e desastrosos da ideia puritana sobre o sexo que
acredito ser imperioso chamar a atenção para a necessidade de se tratar
da questão sexual de modo franco e sem os subterfúgios usualmente
empregados, quando nos preparamos para tratar desse assunto. Em conjunto
com os espíritos mais elevados e livres e o poeta Walt Whitman, eu digo:
``Onde o sexo está faltando, tudo está faltando''. Livremo-nos da
humildade dissimulada tão predominante na superficialidade dos bons
modos. Libertemos o sexo da falsidade e da degradação e percebamos logo
de uma vez por todas que o sexo é um fator de importância central para a
saúde e a harmonia na vida e na arte.
